% ============================================================
% 等离子体物理自学教材
% ============================================================
\documentclass[12pt,a4paper,openany,oneside]{book}
\usepackage[UTF8]{ctex}
\usepackage{geometry}
% ---- 版面密度规格 v1.9.0（2026-10-11）------------------------------------
% 目标：不缩小字号（正文仍 12pt）的前提下，让每页承载更多内容，并重新规划
% 页面元素的垂直间距。边距与间距的实测效果见 processed.md 的 v1.9.0 一节。
% 边距：left 2.1->1.75cm、right 2.6->1.95cm、top 2.1->1.85cm、bottom 1.8->1.55cm。
%   版心由 16.30cm x 25.80cm 增至 17.30cm x 26.30cm，可用面积 +8.1%。
%   行宽 17.30cm 在 12pt 中文下约 41 字/行，仍在可读范围内（不超过 45 字）。
% right 收到 1.95cm 的约束：难度边注盒（marginparwidth + marginparsep）右边界
%   须落在纸面 21.0cm 内。现取 1.1 + 0.15 = 1.25cm，右边界
%   1.75+17.30+0.15+1.10 = 20.30cm，距裁切边 0.70cm。边注内容实测最宽为
%   「[基础]」「[提高]」「[应用]」等 2 字标签（footnotesize 约 0.81cm），
%   在 1.10cm 的 marginparwidth 内不会被拆行或溢出。
% headsep 0.5->0.4cm、footskip 0.8->0.7cm：页眉页脚同步贴近版心。
\geometry{left=1.75cm,right=1.95cm,top=1.85cm,bottom=1.55cm,
          footskip=0.7cm,headsep=0.4cm}
\setlength{\headheight}{15pt}
\usepackage{amsmath,amssymb,amsthm}
\usepackage{mathtools}
\usepackage{bm}
\usepackage{esint}
\usepackage{siunitx}
\usepackage{tikz,pgfplots}
\pgfplotsset{compat=1.18}
\usetikzlibrary{arrows.meta,decorations.pathmorphing,decorations.markings,
  patterns,calc,positioning,shapes,shadings,fadings,through}
\pgfdeclarelayer{bg}
\pgfdeclarelayer{fg}
\pgfsetlayers{bg,main,fg}
\usepackage{xcolor}
\usepackage{graphicx}
\usepackage{enumitem}
% ---- 列表全局规范：替代全书零散的逐实例配置 ----
\setlist{nosep,leftmargin=1.25em,topsep=1.4pt plus 0.5pt minus 0.5pt,
         itemsep=0.5pt plus 0.3pt,parsep=0pt,partopsep=0pt}
\setlist[itemize,1]{label=\textbullet}
\setlist[itemize,2]{label=$\circ$}
\setlist[enumerate,1]{label=\arabic*.}
\setlist[enumerate,2]{label=(\alph*)}
\usepackage{etoolbox}
\usepackage{microtype}
\usepackage{fancyhdr}
\usepackage{titlesec}
\usepackage{tcolorbox}
\tcbuselibrary{theorems,skins,breakable}
\usepackage{booktabs}
\usepackage{array}
\usepackage{longtable}
\usepackage{caption}
% ---- 封面与全书主题色 ----
\definecolor{coverDeep}{RGB}{6, 10, 32}        % 深空藏青（渐变锚点）
\definecolor{coverMid}{RGB}{28, 42, 96}        % 中段藏青蓝（全文标题主色）
\definecolor{coverGlow}{RGB}{196, 92, 42}      % 太阳风暖辉
\definecolor{imfLine}{RGB}{120, 190, 235}      % IMF 磁力线青
\definecolor{swFlow}{RGB}{255, 186, 96}        % 太阳风流线琥珀
\definecolor{magnetopause}{RGB}{120, 232, 190} % 磁层顶青绿
\definecolor{reconX}{RGB}{255, 238, 140}       % X 点高亮
\definecolor{earthBlue}{RGB}{42, 104, 176}     % 地球盘面
% ---- 封面示意图配色（语义化命名，便于逐元素调整）----
\colorlet{shockC}{magnetopause}                % 弓形激波（虚线）
\colorlet{mpC}{magnetopause}                   % 磁层顶
\colorlet{tailC}{imfLine}                      % 磁尾尾瓣磁场
\colorlet{lineC}{imfLine}                      % 磁力线（同 IMF 青）
\colorlet{cuspC}{swFlow}                       % 极尖区标记
\definecolor{plasmaC}{RGB}{46, 40, 96}         % 等离子体片（深紫）
\definecolor{neutralC}{RGB}{10, 10, 20}        % 中性片（近黑）
\definecolor{bLineC}{RGB}{120, 190, 235}       % 封面磁力线
\definecolor{orbC}{RGB}{255, 186, 96}         % 封面回旋轨迹
\definecolor{torusC}{RGB}{96, 168, 224}       % 封面环向磁场
\definecolor{glowC}{RGB}{86, 196, 236}        % 封面电离辉光
\definecolor{filC}{RGB}{150, 226, 255}        % 封面放电丝
\colorlet{sheetC}{swFlow}                      % 等离子体片 / 电流片
\colorlet{earthC}{earthBlue}                   % 地球盘面
% ---- 封面柔光淡出（必须声明在导言区，放进 picture 会扰乱 overlay 定位）----
\tikzfading[name=glowFade, inner color=transparent!85, outer color=transparent!0,
            middle color=transparent!35]
\tikzfading[name=warmFade, inner color=transparent!55, outer color=transparent!0,
            middle color=transparent!18]
\tikzfading[name=xFade, inner color=transparent!75, outer color=transparent!0,
            middle color=transparent!25]
% ---- 图表标题：表格与图注分区处理 ----
\captionsetup{font=small,labelfont={bf,color=coverMid},labelsep=quad,
              justification=centering,skip=3pt,belowskip=1pt}
\captionsetup[table]{position=top,skip=2pt}
\captionsetup[figure]{position=bottom}
% ---- 表格字号与列距统一（不改 arraystretch，避免覆盖 4312 行的显式设置）----
\AtBeginEnvironment{tabular}{\setlength{\tabcolsep}{4pt}\small}
\usepackage{subcaption}
\usepackage[pdfusetitle,
  pdftitle={等离子体物理自学教材·配套手册 v1.32.33},
  pdfauthor={Han \& Kimi},
  pdfsubject={经典弱耦合等离子体物理入门},
  pdfkeywords={等离子体, 聚变, 低温等离子体, 诊断, 自学教材},
  pdflang={zh-CN},
  pdfcreationdate={D:20260913000000Z}]{hyperref}
\hypersetup{colorlinks=true,linkcolor=coverMid,citecolor=coverMid,urlcolor=coverMid,pdfdisplaydoctitle=true}
\usepackage{cleveref}
\usepackage{braket}
\usepackage{float}
\usepackage{multirow}
\usepackage{marginnote}
\usepackage{epigraph}
\usepackage{pifont}          % \ding：圈号 （1）（2）（3） 的可靠字形来源（lmroman 无 U+2460..U+2462）
\usepackage{fvextra}         % verbatim 可断行，消除附录 G 代码越界
\setlength{\marginparwidth}{1.1cm}
\setlength{\marginparsep}{0.15cm}
% ---- 圈号 （1）（2）（3） 的安全写法 ----
% 原实现直接在 lmroman 中排 U+2460..U+2462，该字体无此字形，PDF 中渲染为空白。
% \ding{172..174} 对应 （1）..（3），pifont 的 dingbat 字体覆盖完整。
\newcommand{\cnum}[1]{\ding{\numexpr171+#1\relax}}
% ---- 段落、公式与浮动体间距：回收造成大片留白的垂直空间 ----
% v1.27.1：在不改字号、不改页边距、不删一个字的前提下再收一档行间空白，
% 把新增实拍图占掉的页数压回来（段落靠首行缩进区分，不靠空行）。
\setlength{\parskip}{1.0pt plus 0.4pt minus 0.2pt}
\setlength{\jot}{1.0pt}
\AtBeginDocument{%
  \setlength{\abovedisplayskip}{2.2pt plus 0.8pt minus 0.8pt}%
  \setlength{\belowdisplayskip}{2.2pt plus 0.8pt minus 0.8pt}%
  \setlength{\abovedisplayshortskip}{1.8pt plus 0.6pt}%
  \setlength{\belowdisplayshortskip}{1.8pt plus 0.6pt}%
}
% 浮动体上下间距（默认 book 类为 12pt 级，是页面上最显眼的空白来源）
\setlength{\textfloatsep}{5pt plus 2pt minus 2pt}
\setlength{\floatsep}{5pt plus 2pt minus 2pt}
\setlength{\intextsep}{5pt plus 2pt minus 2pt}
% 放宽浮动体在同一页的比例上限，减少"图表被挤到下一页"留下的整块空白
\renewcommand{\topfraction}{0.92}
\renewcommand{\bottomfraction}{0.75}
\renewcommand{\textfraction}{0.06}
\renewcommand{\floatpagefraction}{0.72}
\setcounter{topnumber}{3}
\setcounter{bottomnumber}{2}
\setcounter{totalnumber}{4}
\pagestyle{fancy}
\fancyhf{}
% 中文页眉不用 \itshape：SimSun 无真斜体，会触发合成倾斜，与正文正体断裂。
\fancyhead[L]{\nouppercase{\small\heiti\leftmark}}
\fancyhead[R]{\small\bfseries\thepage}
\renewcommand{\headrulewidth}{0.3pt}
% 章首页（\chapter 下发 \thispagestyle{plain}）原先无页眉也无页码，此处补上页码。
\fancypagestyle{plain}{%
  \fancyhf{}%
  \fancyhead[R]{\small\bfseries\thepage}%
  \renewcommand{\headrulewidth}{0.3pt}%
}
% 附录章：\chapter* 不更新 \leftmark，导致 186--210 页页眉恒为"第 12 章"。
% \appchapter 显式维护页眉与书签，并让小节/定理编号继承附录章号（A/B/C...）。
% v1.25.0：字母方案 A--Z, AA, AB…。\thechapter 必须在这里就 \xdef 成**纯文本**：
% 若留成含 \ifnum / expl3 的宏，hyperref 生成书签（PDF 字符串）时会展开失败，
% 书签会丢掉「AA.」前缀，pdf_outline 的 O3a 就会报「印刷目录与书签不同名」。
\makeatletter
\newcommand{\appendixAlph}[1]{%
  \ifnum#1>26 A\@Alph{\numexpr#1-26\relax}\else\@Alph{#1}\fi}
\makeatother
\newcommand{\appchapter}[1]{%
  \cleardoublepage
  \refstepcounter{chapter}%
  \xdef\thechapter{\appendixAlph{\arabic{chapter}}}%
  \chapter*{#1}%
  \addcontentsline{toc}{chapter}{#1}%
  \markboth{#1}{#1}%
  \thispagestyle{plain}%
}
% 前言/符号说明/结语等既非正文章、也不属于附录的部分：
% 需要页眉与书签，但绝不能递增 chapter 计数器（否则正文章号整体后移）。
\newcommand{\frontchapter}[1]{%
  \cleardoublepage
  \chapter*{#1}%
  \addcontentsline{toc}{chapter}{#1}%
  \markboth{#1}{#1}%
  \thispagestyle{plain}%
}
\titleformat{\chapter}[display]
  {\normalfont\huge\heiti\color{coverMid}}
  {\filleft\fontsize{52}{52}\selectfont\color{coverMid!30}\thechapter}
  {-2.4ex}
  {\Huge\heiti\filright}
  [\vspace{0.35ex}{\color{coverMid!45}\titlerule[1.2pt]}]
\titlespacing*{\chapter}{0pt}{-30pt}{15pt}
% ---- 难度边注：就地排版，不用 \hfill 顶边 ----
% 历史缺陷：
%  (a) 旧宏内 \hfill 使 2pt \rule 紧贴边注盒右边界，墨迹到达 20.62cm，越出盒边界
%      (20.50cm)，影响 20 页；配合 geometry right 2.3cm->2.6cm 后已落回盒内。
%  (b) 旧宏锚定的是"当前位置"，写在 \section{...} 之后时会掉到标题下方的正文
%      首行。现改为在 \section 的"标题后"钩子里立刻排出，锚点即标题行。
\newcommand{\diffmarktext}[1]{%
  \textcolor{coverMid!80}{\rule[-0.1em]{2pt}{0.9em}}\,%
  \textcolor{coverMid}{\footnotesize\textbf{[#1]}}%
}
% 就地排出：\section{...}\difficulty{基础} 中该宏紧跟标题文本，仍在标题行内，
% \marginnote 因此锚定标题行本身。不再依赖额外钩子，也不会漏排。
\newcommand{\difficulty}[1]{\marginnote{\raggedright\diffmarktext{#1}}}
% 兼容旧用法别名（tcolorbox 标题内等场合），行为与 \difficulty 一致
\newcommand{\difficultyhere}[1]{\marginnote{\raggedright\diffmarktext{#1}}}
\titleformat{\section}
  {\normalfont\Large\heiti\color{coverMid}}{\thesection}{0.9em}{}
\titlespacing*{\section}{0pt}{7pt plus 1.6pt minus 1.6pt}{3.4pt plus 1.2pt minus 0.8pt}
\titleformat{\subsection}
  {\normalfont\large\heiti\color{coverMid!85}}{\thesubsection}{0.8em}{}
\titlespacing*{\subsection}{0pt}{5.4pt plus 1.2pt minus 0.8pt}{2.4pt}
% ---- 章前引语：楷体，与标题黑体形成字型区分（KaiTi 无斜体，不设 itshape）----
\setlength{\epigraphwidth}{0.62\textwidth}
\renewcommand{\epigraphsize}{\small\kaishu}
\setlength{\epigraphrule}{0pt}
\definecolor{thmblue}{RGB}{41, 98, 255}
\definecolor{defgreen}{RGB}{46, 125, 50}
\definecolor{cororange}{RGB}{230, 81, 0}
\definecolor{exampurple}{RGB}{106, 27, 154}
\definecolor{insightred}{RGB}{183, 28, 28}
\definecolor{warnyellow}{RGB}{255, 160, 0}
\definecolor{historybrown}{RGB}{121, 85, 72}
\definecolor{reviewblue}{RGB}{25, 118, 210}
\definecolor{tipgreen}{RGB}{56, 142, 60}
% ---- 插图统一视觉语言（2026-10-19 按 docs/palette.md 的「深蓝 + 金」重锚）----
% 同一物理量在全书各图里用同一色：
%   电子/快/束流 = figGold（亮金，只作线与填充）；离子/慢/背景 = figBlue（钢蓝）；
%   磁场/磁力线  = figGreen（深青）；判据/阈值 = figVerm（赭红）；第三分类 = figPurple；
%   辅助线/包络 = figGray（灰）虚线；金色文字一律用 figAmber（暗金）。
% 线宽：坐标轴 0.9pt、数据曲线 1.2pt、辅助虚线 0.6pt；箭头统一 Stealth。
% 字号：图内主标签 \small、次标签与刻度 \footnotesize；内容标签全书禁用极小字号。
% 标签永不压线：必须靠近曲线时用 fill=white, inner sep=1.5pt 的光环；
% 区域标签若放不进色块（留不出约 3mm 边）就移到块外，用 figlead 引线点指。
% 图例一律放在绘图区外（下方或右侧），带线型样本。
\definecolor{figInk}{HTML}{12203A}
\definecolor{figNavy}{HTML}{1B2E52}
\definecolor{figBlue}{HTML}{2F5286}
\definecolor{figSky}{HTML}{6C86AE}
\definecolor{figGreen}{HTML}{2A7373}
\definecolor{figGold}{HTML}{C2A14D}
\definecolor{figAmber}{HTML}{8A6520}
\definecolor{figVerm}{HTML}{9C443A}
\definecolor{figPurple}{HTML}{5F4066}
\definecolor{figGray}{HTML}{6B7280}
\definecolor{figPaper}{HTML}{F6F2E7}
\colorlet{figOrange}{figGold}
\colorlet{figGoldDeep}{figAmber}
\tikzset{
  figaxis/.style={draw=figInk, line width=0.9pt, ->, >=Stealth},
  figcurve/.style={line width=1.2pt},
  fighelp/.style={draw=figGray, line width=0.6pt, dashed},
  fignode/.style={rounded corners=2pt, line width=0.8pt, inner sep=4pt},
  figlead/.style={draw=figGray, line width=0.5pt, ->, >=Stealth, shorten >=1.5pt},
  figtick/.style={font=\footnotesize, text=figInk},
}
% ---- 真实图像配图（2026-10-16 起）：开放许可照片/机构示意图的排版约定 ----
% 只用 CC0 / 公有领域 / CC BY / CC BY-SA 的图像；每张实拍图随图署名，
% 在 \caption 之后用 \figsource{作者｜许可｜来源} 写一行灰色小字（作者与许可以 CREDITS.md 为准）。
% 排版宽度：单栏整幅不超过 0.86\textwidth，双图并排各 0.42\textwidth；
% 正文里的图一律 figure[htbp] + \centering，图注首句点题（这一条从 v1.27.0 起用于新增配图）。
\newcommand{\figsource}[1]{\par\vspace{1pt}{\footnotesize\color{figGray}图片来源：#1}}
\newcommand{\predict}[1]{\par\smallskip{\footnotesize\color{figGray}\textbf{先猜再看}：#1}\par\smallskip}
\newtcbtheorem[number within=chapter]{theorem}{定理}{
  colback=thmblue!5,colframe=thmblue,fonttitle=\bfseries,breakable,
  left=4pt,right=4pt,top=4pt,bottom=4pt}{thm}
\newtcbtheorem[number within=chapter]{definition}{定义}{
  colback=defgreen!5,colframe=defgreen,fonttitle=\bfseries,breakable,
  left=4pt,right=4pt,top=4pt,bottom=4pt}{def}
\newtcbtheorem[number within=chapter]{corollary}{推论}{
  colback=cororange!5,colframe=cororange,colbacktitle=cororange!62!black,fonttitle=\bfseries,breakable,
  left=4pt,right=4pt,top=4pt,bottom=4pt}{cor}
\newtcbtheorem[number within=chapter]{example}{例题}{
  colback=exampurple!5,colframe=exampurple,fonttitle=\bfseries,breakable,
  left=4pt,right=4pt,top=4pt,bottom=4pt}{ex}
\newtcbtheorem[number within=chapter]{insight}{物理洞见}{
  colback=insightred!5,colframe=insightred!80,fonttitle=\bfseries,breakable,
  left=4pt,right=4pt,top=4pt,bottom=4pt}{ins}
\newtcbtheorem[number within=chapter]{warning}{常见误区}{
  colback=warnyellow!8,colframe=warnyellow,colbacktitle=warnyellow!55!black,fonttitle=\bfseries,breakable,
  left=4pt,right=4pt,top=4pt,bottom=4pt}{warn}
\newtcbtheorem[number within=chapter]{history}{物理史话}{
  colback=historybrown!5,colframe=historybrown,fonttitle=\bfseries,breakable,
  left=4pt,right=4pt,top=4pt,bottom=4pt}{hist}
\newtcbtheorem[number within=chapter]{review}{本节要点}{
  colback=reviewblue!5,colframe=reviewblue,fonttitle=\bfseries,breakable,
  left=4pt,right=4pt,top=4pt,bottom=4pt}{rev}
\newtcbtheorem[number within=chapter]{tip}{学习提示}{
  colback=tipgreen!5,colframe=tipgreen,colbacktitle=tipgreen!62!black,fonttitle=\bfseries,breakable,
  left=4pt,right=4pt,top=4pt,bottom=4pt}{tip}
\newtcolorbox{miniquiz}[1][]{
  colback=reviewblue!3,colframe=reviewblue,fonttitle=\bfseries,
  title={随堂自测},breakable,
  left=4pt,right=4pt,top=4pt,bottom=4pt,#1}
% 随堂自测的参考要点（2026-10-16 起）：先自己答，再看这一行
\newcommand{\quizans}[1]{\par\medskip{\color{figGray}\rule{\linewidth}{0.3pt}}\par\smallskip
  {\footnotesize\textbf{参考要点：}#1}}
% 章末概念图（2026-10-16 起，借鉴 Novak & Gowin 的概念图与 OpenStax 的「概念小结」）
\tikzset{
  cmnode/.style={draw,rounded corners=2pt,align=center,inner sep=3pt,font=\footnotesize,minimum height=7mm},
  cmedge/.style={-{Latex[length=1.6mm]},draw=figGray,line width=0.7pt},
  cmA/.style={cmnode,draw=figBlue,fill=figBlue!8},
  cmB/.style={cmnode,draw=figOrange,fill=figOrange!10},
  cmC/.style={cmnode,draw=figGreen,fill=figGreen!10},
  cmD/.style={cmnode,draw=figVerm,fill=figVerm!8},
  cmE/.style={cmnode,draw=figGray,fill=figGray!8},
  rmnode/.style={cmnode,fill=white,minimum width=25mm},
  rmedge/.style={-{Latex[length=1.6mm]},draw=figGray,line width=0.7pt},
  rmABC/.style={rmnode,draw=figGreen,very thick},
  rmAB/.style={rmnode,draw=figBlue,very thick},
  rmBC/.style={rmnode,draw=figOrange,very thick},
  nxchip/.style={draw=figInk!70,rounded corners=1.5pt,minimum width=1.05cm,minimum height=0.95cm,font=\footnotesize,inner sep=1pt},
}
\newtcolorbox{chapterreview}[1][]{
  colback=tipgreen!5,colframe=tipgreen,fonttitle=\bfseries,
  title={本章要点回顾},breakable,
  left=4pt,right=4pt,top=4pt,bottom=4pt,#1}
\newtcolorbox{formula}[1][]{
  colback=blue!5,colframe=blue!60,fonttitle=\bfseries,
  title={核心公式速查},breakable,
  left=4pt,right=4pt,top=4pt,bottom=4pt,#1}
% 兼容两参数签名的速查卡（{标题}{key}）：key 不再掉入正文
\NewTColorBox{formulabox}{m m}{%
  colback=blue!5,colframe=blue!60,fonttitle=\bfseries,breakable,
  title={核心公式速查：#1},
  phantomlabel={frml:#2},
  left=4pt,right=4pt,top=4pt,bottom=4pt}
\newtcolorbox{checklist}[1][]{
  colback=green!5,colframe=green!60,fonttitle=\bfseries,
  title={自学检查清单},breakable,
  left=4pt,right=4pt,top=4pt,bottom=4pt,#1}
% 兼容两参数签名的检查清单（{标题}{key}）：key 不再掉入正文
\NewTColorBox{checklistbox}{m m}{%
  colback=green!5,colframe=green!60,fonttitle=\bfseries,breakable,
  title={#1},
  phantomlabel={ckl:#2},
  left=4pt,right=4pt,top=4pt,bottom=4pt}
\newtcolorbox{strategy}[1][]{
  colback=figOrange!5,colframe=figOrange,fonttitle=\bfseries,
  title={解题策略},breakable,
  left=4pt,right=4pt,top=4pt,bottom=4pt,#1}
% ---- 统一色块样式：必须在上面 16 个定义之后，否则被各自的 left=4pt 反压 ----
% v1.9.0：内边距 5/3pt -> 4.5/2.5pt，块前后间距 5pt -> 3.5pt，标题上下 2pt -> 1.5pt。
\tcbset{
  boxsep=0pt,
  left=4.5pt,right=4.5pt,top=2.5pt,bottom=2.5pt,
  before skip=3.5pt plus 1.5pt minus 1pt,
  after skip=3.5pt plus 1.5pt minus 1pt,
  fonttitle=\heiti\small,
  toptitle=1.5pt,bottomtitle=1.5pt,
  titlerule=0.4pt,
  sharp corners=southwest,
  enhanced
}
\newcommand{\diff}{\mathrm{d}}
\newcommand{\ii}{\mathrm{i}}
\newcommand{\ee}{\mathrm{e}}
\newcommand{\eps}{\varepsilon}
\newcommand{\vect}[1]{\bm{#1}}
\newcommand{\unitvec}[1]{\hat{\vect{#1}}}
\newcommand{\grad}{\nabla}
\newcommand{\curl}{\nabla\times}
\newcommand{\divg}{\nabla\cdot}
\newcommand{\lap}{\nabla^2}
\newcommand{\tens}[1]{\bm{#1}}
\newenvironment{chapterintro}{%
  \begin{quotation}\itshape\small\setlength{\topsep}{2pt}\setlength{\parskip}{1.2pt}%
}{%
  \end{quotation}%
}
% 附录编号进入两位数后（AA...AM），目录 \numberline 的 2.3em 数字列不够宽，
% 每个附录小节目录行会溢出约 1.1 pt；把数字列加宽到 2.5em（缩进与正文不变）。
\makeatletter
\renewcommand*\l@section{\@dottedtocline{1}{1.5em}{2.5em}}
\makeatother
\begin{document}

% ======================= 封面 =======================
\thispagestyle{empty}
\begin{tikzpicture}[remember picture, overlay]
  \fill[coverDeep] (current page.north west) rectangle (current page.south east);
  \begin{pgfonlayer}{bg}
    \fill[path fading=warmFade, color=coverGlow!55!white]
      ([xshift=-6.2cm,yshift=3.2cm]current page.north west) circle (13cm);
    \fill[path fading=warmFade, color=coverGlow!32!white]
      ([xshift=-7.0cm,yshift=4.2cm]current page.north west) circle (8cm);
    \fill[path fading=glowFade, color=black]
      ([xshift=1.2cm,yshift=-1.2cm]current page.south east) circle (14cm);
  \end{pgfonlayer}
  % 三条横线：示意「学习法—书单—题库」三件事
  \fill[white!22] ([xshift=2.6cm,yshift=-3.7cm]current page.north west) rectangle ++(9.6cm,0.9pt);
  \fill[white!16] ([xshift=2.6cm,yshift=-4.5cm]current page.north west) rectangle ++(7.4cm,0.9pt);
  \fill[white!10] ([xshift=2.6cm,yshift=-5.3cm]current page.north west) rectangle ++(5.2cm,0.9pt);
\end{tikzpicture}

\vspace*{2.6cm}
\begin{flushleft}
\hspace*{1.6cm}{\color{white}\heiti\fontsize{16}{20}\selectfont 等离子体物理自学教材\quad 配套手册}\\[1.4em]
\hspace*{1.6cm}{\color{white}\heiti\fontsize{40}{46}\selectfont 从学习法到题库}\\[0.8em]
\hspace*{1.6cm}{\color{white!82!coverGlow}\heiti\fontsize{20}{26}\selectfont 学习策略 · 经典教材逐章对照 · 跨教材例题 · 分章自测}\\[2.6em]
\hspace*{1.6cm}{\color{white!78}\large 主册（12 章正文 + 48 个专题附录）的配套读物}\\[0.5em]
\hspace*{1.6cm}{\color{white!78}\large 主册附录 I、J、K、L 的完整内容，单独成册，便于边读边查}
\end{flushleft}

\vfill
\hspace*{1.6cm}{\color{white!72}\large Han \& Kimi}\hfill
{\color{white!72}\large v1.32.33}\hfill
{\color{white!60}\large 2026 年 10 月}\hspace*{1.6cm}
\vspace*{1.2cm}

\clearpage
\thispagestyle{empty}
\frontchapter{这本手册怎么用}

四篇各管一件事：

\begin{itemize}
\item \textbf{附录 I（学习策略与刻意练习手册）}：怎么读、怎么练、怎么复习。合书复述、间隔重复、交错练习、错题本、十二周节奏、读论文三遍法，都在这一篇里。建议在读正文之前先翻一遍。
\item \textbf{附录 J（经典教材逐章对照与延伸阅读）}：本书每一章对应 Chen、Fitzpatrick 等经典教材的哪一章。想继续深入时该读哪本专著，也在这里查。
\item \textbf{附录 K（跨教材经典例题精选）}：八道跨教材例题，覆盖德拜屏蔽、磁镜、漂移、Alfv\'en 波、Langmuir 探针。每题都有完整解答。这是正文例题的补充。
\item \textbf{附录 L（分章自测题库）}：12 章 $\times$ 6 题的题库（每题都带答案），按章刷题时用它。
\end{itemize}

\textbf{为什么单独成册。}查书单时读者不在读正文，刷题时也不想被推导带走，这四篇便单收一本。主册因此更薄，能一口气读完；要查时再翻开这本手册。

\textbf{字母为什么跳过 I--L。}主册的附录字母没有重排：I、J、K、L 的位置留给这本手册。主册附录 H 之后直接接 M，两边字母一一对应，引用时不会错位。

\textbf{版本与许可。}本手册与主册同版本（v1.32.33），许可同为 CC BY-NC-SA 4.0。源码与四篇 \texttt{.tex} 都在仓库 \texttt{src/} 下；正文引用的脚本在 \texttt{tools/calc/} 与 \texttt{code/} 下。

\clearpage
\tableofcontents

\clearpage
\setcounter{chapter}{8}

\appchapter{附录 I：学习策略与刻意练习手册}
%======================================================================
% 附录I：学习策略与刻意练习手册
%======================================================================
% 章标题 \appchapter{附录I：...} 由主文件 plasma_physics_textbook_v2.tex 发出
% （检查套件按主文件判定附录字母连续性），本文件只提供正文。

前面八个附录回答的是"等离子体物理是什么"，这一个附录回答的是"怎么把它学进脑子里"。

本书是自学教材：没有老师点名，没有考试倒逼，也没有人替你判断"到底懂没懂"。自学最大的风险通常不是内容太难，而是\textbf{用低效的方法把时间花光}——花了三个月，书翻旧了，合上书却推不出德拜长度。

本附录把学习科学中经过实验检验的做法，翻译成可以直接执行的动作，并逐条指出它在本书里的落点。

\begin{insight}{本附录与前言的分工}{}
前言里的"三类自学人群路线图"回答\textbf{先读什么、读多快、跳过什么}；本附录回答\textbf{读完之后怎么巩固、隔多久复习、怎么证明自己真的会了}。两者合起来才是一条完整的学习闭环：前言负责排课，本附录负责验收。
\end{insight}


\section{三个把"读懂了"误当"学会了"的模式}

\begin{warning}{流畅性错觉}{}
重新阅读、划线、抄公式，都会让材料\textbf{看起来}越来越熟悉。这种熟悉感来自"再认"（看到答案能认出来），而不是"回忆"（不给提示能写出来）。考试与科研需要的是后者，而前者会稳定地高估自己的掌握程度。
\end{warning}

\textbf{误区一：重读与划线。}
把一章读三遍，比读一遍再合书复述一遍，效果差得多：前者提高的是对排版的熟悉度，后者才是在建立提取路径。
诊断方法很简单：读完一节后立刻合上书，在白纸上写出这一节的\textbf{核心方程、适用条件、一个反例}。写不出来，说明刚才那一遍只是"看过了"。

\textbf{误区二：集中突击。}
连续 6 小时攻同一章，短期感觉很爽；两周后回看，遗忘得比"每天 1 小时、分 6 天"更彻底。
同样总时长下，把学习\textbf{分散}到多天，长期保持率显著更高——这条结论在跨度几十年的元分析中反复出现（见 I.7 文献表）。

\textbf{误区三：同类题目连做。}
把 10 道双流不稳定性题目连着做完，正确率会虚高，因为每道题都在重复同一个套路；一旦题目混在一起、需要自己判断"该用哪个方法"，正确率立刻掉下来。
更有效的做法是\textbf{交错练习}：把第 6 章的波、第 7 章的动理学、第 9 章的不稳定性题目打乱混做，逼迫自己在解题前先做一次"方法选择"。

本书中，第 9 章补充材料的不稳定性分类树、附录 F 的 23 条高频易错点，都是为这种"先判断再计算"的训练准备的。


\section{五个可执行动作}

\subsection{动作一：合书复述（检索练习）}

\textbf{怎么做：}读完一节，立刻合书，用 3 分钟在白纸上写出：核心方程 $\to$ 每个符号的物理含义 $\to$ 推导的关键一步 $\to$ 成立条件 $\to$ 一个会被它证伪的反例。写完再翻书核对，\textbf{用红笔标出漏掉的那一条}。

\textbf{在本书的落点：}每章的"本章速查卡"（补充材料末尾的 `formulabox`）就是现成的核对清单；每章末尾的"自学检查清单"（`checklistbox`）可直接当复述提纲。

\textbf{常见走样：}变成"抄笔记"。抄写时眼睛一直盯着书，提取路径没有被使用，效果接近于零。

\subsection{动作二：间隔重复（1 / 3 / 7 / 21 天）}

\textbf{怎么做：}同一节内容，在学完后的第 1、3、7、21 天各做一次 2--3 分钟的合书复述。四次都过关，可以降为每月一次。

\textbf{为什么是这四个点：}遗忘在最初几天最快。复习的价值不在于"再看一遍"，而在于\textbf{在快要忘掉的时候把它取回来}——取回越费力（在成功的前提下），记忆越牢。

\textbf{在本书的落点：}把每章速查卡抄成 12 张卡片（一章一张），在卡片背面只写标题与日期，正面留白。复习时只看日期与标题，凭记忆默写符号表。

\subsection{动作三：交错练习}

\textbf{怎么做：}每周留出一次"混合练习"：从本周学习内容与之前 2--3 章中各抽 3 道题，打乱顺序限时完成。重点是\textbf{在动笔前用一句话写出"我为什么选这个方法"}。

\textbf{在本书的落点：}附录 K 的八道跨教材例题覆盖德拜屏蔽、磁镜、漂移、Alfvén 波、Landau 阻尼、Spitzer 电阻率、鞘层与重联，天然适合交错使用；把它们与各章例题混排出题效果最好。

\subsection{动作四：自我解释}

\textbf{怎么做：}每推完一步问自己三句话——"这一步用了什么假设？""如果这个假设不成立会怎样？""能不能用一句不含公式的话说出来？"把答案写在一旁。

\textbf{在本书的落点：}本书大量使用\textbf{推导后紧跟假设讨论}的写法（例如第 4 章流体闭合、第 7 章 Landau 阻尼的围道选择）。读到这些段落时不要跳过，它们正是自我解释的范本。

\subsection{动作五：例题读完合书重做}

\textbf{怎么做：}例题不要"读懂"就算过。读完题目后合上书，从空白纸开始独立做一遍；做不出来再看提示，然后\textbf{隔一天再做一遍}。同一道题做三遍（读后、隔日、隔周）比做三道同类新题更有价值。

\textbf{在本书的落点：}第 1--12 章的补充材料各含 5 道经典例题，加上附录 K 的 8 道，共 86 道；建议至少把其中 30 道按"三遍法"完整做过。

\begin{miniquiz}
\textbf{自测：}你的复习计划是"每周末把本周内容重读一遍"。请指出这个计划的两个缺陷，并改写成符合动作一与动作二的版本。（参考：重读不产生提取；一周一次的点太少且过晚。）
\end{miniquiz}


\section{元认知与错题本}

\textbf{错题本不是抄题本。}每道错题只记三栏：

\begin{enumerate}
\item \textbf{错在哪一步}——精确到"把 $\grad p$ 当成了标量""忘了 $\ee$ 的符号""单位换成了 CGS 但公式是 SI"；
\item \textbf{根因归类}——概念混淆（回第几章哪一节）／数学工具（第 1 章）／量纲与量级／审题遗漏；
\item \textbf{下次触发条件}——写成一个可执行的检查动作，例如"凡是出现 $\grad$ 与 $\vect{E}$ 同时出现的式子，先确认量纲"。
\end{enumerate}

三栏里最值钱的是第三栏。它把一次具体错误升级为一条\textbf{可复用的检查动作}，这正是"刻意练习"里"即时反馈 + 针对性重复"的落点：练习必须有明确目标、即时反馈与超出舒适区的难度，单纯重复已经会做的事不叫刻意练习。

\begin{checklistbox}{每周复盘清单}{app-i-weekly}
\begin{itemize}
\item 本周学过的每一节，都做过至少一次\textbf{合书复述}吗？
\item 错题本这周新增了几条？每条都写了"下次触发条件"吗？
\item 有没有做过一次\textbf{混排限时}练习（动作三）？
\item 有没有哪道例题，我在\textbf{不看答案}的情况下完整推出过？
\item 下周要复习的旧章节，是哪几章（按 1/3/7/21 天节奏）？
\end{itemize}
\end{checklistbox}


\section{量级估算与量纲检验}

物理直觉不是玄学，它由一批\textbf{记得住的典型数值}和一套\textbf{三十秒能做完的检查}组成。附录 K 的例题大多可以在两三行内估出结果，请把它们当作"直觉素材"背下来。

\subsection{三个必须记住的量级（氘等离子体）}

\begin{itemize}
\item \textbf{德拜长度}：$T=10\ \mathrm{eV}$、$n=10^{18}\ \mathrm{m^{-3}}$ 时，$\lambda_D\approx 23.5\ \mathrm{\mu m}$，对应 $N_D\approx 5.4\times 10^{4}$ 个粒子——远大于 1，所以流体图像在这里是合法的。
\item \textbf{回旋半径}：$\mathrm{D^+}$ 在 $B=5\ \mathrm{T}$ 下，$T=10\ \mathrm{eV}$ 时 $\rho_L\approx 0.13\ \mathrm{mm}$；$T=1\ \mathrm{keV}$ 时 $\approx 1.3\ \mathrm{mm}$；$T=10\ \mathrm{keV}$ 时 $\approx 4.1\ \mathrm{mm}$。$\rho_L\propto\sqrt{T}$，温度涨 1000 倍只要半径涨约 32 倍。
\item \textbf{Alfvén 速度}：$\mathrm{D^+}$、$n=10^{20}\ \mathrm{m^{-3}}$、$B=5\ \mathrm{T}$ 时 $v_A\approx 7.7\times 10^{6}\ \mathrm{m/s}$，约合光速的 2.6\%；磁场扰动穿过 1 m 只需约 $1.3\times 10^{-7}\ \mathrm{s}$。
\end{itemize}

\subsection{量纲检验三条}

\begin{enumerate}
\item \textbf{逐项同量纲}：方程两边、以及每一项，量纲必须一致。把 $\omega$、$k$、$v_A$、$\lambda_D$ 的量纲列成一张小表，比每次现推更快。
\item \textbf{无量纲组合}：物理结论通常只依赖无量纲量（$k\lambda_D$、$n\lambda_D^3$、$\beta$、$\rho_L/L$、$\omega/\omega_{ci}$）。先问"这个问题的无量纲参数是什么"，往往比直接积分更快看清物理。
\item \textbf{极限行为}：令参数取 $0$ 或 $\infty$ 看结果是否退化成已知情形。Landau 阻尼率在 $k\lambda_D\to 0$ 时应趋于零（$\propto\exp[-1/(2k^2\lambda_D^2)]$ 指数小），若推导结果在这极限下发散，一定错了。
\end{enumerate}

\begin{miniquiz}
\textbf{估算练习：}不查书，估算 $B=3\ \mathrm{T}$、$T_i=5\ \mathrm{keV}$ 的氘离子回旋频率与回旋半径。提示：$\omega_{ci}=qB/m$，$\rho_L=\sqrt{2T/m}/\omega_{ci}$（量级上用 $\sqrt{T/m}$ 即可）。答案应为 $\omega_{ci}\approx 1.4\times 10^{8}\ \mathrm{s^{-1}}$、$\rho_L$ 量级几毫米——如果你算出厘米或微米，请回第 3 章重做量纲。
\end{miniquiz}


\section{十二周节奏模板}

下面的模板假定\textbf{每天 60 分钟、每周 6 天、共 12 周}完成第一遍通读。它与前言的"三类人群路线图"配合使用：前言决定你读哪些章节，本表决定每天这 60 分钟怎么切、周末交什么作业。

\small
\begin{longtable}{@{}p{0.9cm}p{3.1cm}p{6.8cm}p{3.8cm}@{}}
\caption{十二周自学节奏模板（每天 60 分钟 $\times$ 6 天）}\label{tab:app-i-12weeks}\\
\toprule
周次 & 主攻内容 & 每天 60 分钟怎么切 & 周末产出 \\
\midrule
\endfirsthead
\toprule
周次 & 主攻内容 & 每天 60 分钟怎么切 & 周末产出 \\
\midrule
\endhead
\midrule
\multicolumn{4}{r}{\small（续下页）}\\
\endfoot
\bottomrule
\endlastfoot
1 & 第 1 章 数学与统计力学预备知识 & 15 min 复习昨日复述 $\to$ 25 min 读新节 $\to$ 15 min 例题重做 $\to$ 5 min 写速查卡 & 第 1 章清单自评；量纲表一张 \\
2 & 第 2 章 什么是等离子体 & 同上，例题时间加到 20 min & 德拜屏蔽三问（判据/量级/反例） \\
3 & 第 3 章 单粒子运动 & 20 min 读 $\to$ 20 min 画轨道图 $\to$ 20 min 估算 $\rho_L$、$\omega_c$ & 三类漂移方向默写；磁镜例题重做 \\
4 & 第 4 章 流体描述基础 & 25 min 推导 $\to$ 20 min 自我解释假设 $\to$ 15 min 速查卡 & 双流体方程闭合条件复述 \\
5 & 第 5 章 磁流体力学 & 25 min 读 $\to$ 20 min 例题 $\to$ 15 min 交错练习（混入第 4 章） & MHD 波三支速度量级估算 \\
6 & 第 6 章 等离子体中的波 & 20 min 读 $\to$ 25 min 画色散曲线 $\to$ 15 min 极限检验 & 冷等离子体色散关系默写 \\
7 & 第 7 章 动理学理论 & 25 min 读 $\to$ 20 min 分布函数练习 $\to$ 15 min 复述 Landau 围道 & Landau 阻尼率量级估算 \\
8 & 第 8 章 碰撞与输运 & 20 min 读 $\to$ 25 min 输运系数估算 $\to$ 15 min 错题本 & Spitzer 电阻率与扩散时间估算 \\
9 & 第 9 章 等离子体不稳定性 & 20 min 读 $\to$ 25 min 判据练习 $\to$ 15 min 分类树默画 & 不稳定性分类树；三种不稳定性判据对比 \\
10 & 第 10 章 约束与受控核聚变 & 25 min 读 $\to$ 20 min 劳森判据估算 $\to$ 15 min 交错练习 & 劳森判据三参数扫描手算 \\
11 & 第 11 章 低温工业等离子体 & 20 min 读 $\to$ 25 min 鞘层估算 $\to$ 15 min 复述玻姆判据 & 鞘层厚度与离子通量估算 \\
12 & 第 12 章 诊断 + 全局复习 & 30 min 诊断原理 $\to$ 30 min 附录 F/G/H 串读 & 全书 23 条易错点自查；附录 H 符号表抽测 \\
\end{longtable}
\normalsize

\textbf{进度落后怎么办：}宁可推迟一周，也不要把"合书复述"和"例题重做"砍掉——那两项才是模板的全部价值所在。可以先缓读第 11、12 章等相对独立的章节，但第 1--9 章的概念依赖链不能跳。


\section{自测题四型}

自测题的质量决定复习的质量。下面四型覆盖了从记忆到迁移的完整梯度，每型给出本书中的一个示例与判分标准。

\begin{enumerate}
\item \textbf{判据型——"什么时候可以用这个方法？"}\\
例：什么时候可以把等离子体当流体处理，什么时候必须用动理学？\\
判分：必须同时说出\textbf{判据}（$N_D\gg 1$、$k\lambda_D\ll 1$、碰撞足够频繁使分布接近 Maxwell 分布）与\textbf{失效应}（无碰撞、波--粒子共振时流体图像失效）。

\item \textbf{量级型——"答案大概是多少？"}\\
例：$B=5\ \mathrm{T}$、$T=10\ \mathrm{eV}$ 的氘离子回旋半径是多少？\\
判分：答案应落在 $0.1$--$0.2\ \mathrm{mm}$；数量级对即算过，但必须能说出 $\rho_L\propto\sqrt{T}/B$ 的来源。

\item \textbf{反例型——"这句话反着说对不对？"}\\
例："温度升高，碰撞频率一定升高。"请给出反例。\\
判分：指出库仑碰撞频率 $\nu\propto nT^{-3/2}$ 随温度\textbf{下降}，而与中性粒子的碰撞才随温度升高；能区分两种碰撞机制即算过。

\item \textbf{迁移型——"把方法搬到新问题上。"}\\
例：磁镜损失锥的几何推理（$\sin^2\theta_c=1/R$）能否搬到环形约束装置的俘获粒子分析上？\\
判分：能指出哪一部分可以搬（绝热不变量与速度空间判据）、哪一部分不能（磁场几何与 $B_{\max}/B_{\min}$ 的构造方式不同）。
\end{enumerate}

建议每章学完后，自己按四型各出一道题并交换给同伴（或自我隔周作答）。\textbf{能出题，才说明你看清了这一章的骨架。}


\section{出处与延伸阅读}

本附录的主张来自学习科学文献与经典解题方法论，列在下面以便读者自行核对。物理教材的对照关系见附录 J，分级书单见附录 E。

\normalsize
{\small
\begin{longtable}{@{}p{4.4cm}p{5.2cm}p{5.2cm}@{}}
\caption{本附录引用的学习科学文献与出处}\label{tab:app-i-refs}\\
\toprule
文献 & 出处 & 本附录中的用途 \\
\midrule
\endfirsthead
\toprule
文献 & 出处 & 本附录中的用途 \\
\midrule
\endhead
\midrule
\multicolumn{3}{r}{\small（续下页）}\\
\endfoot
\bottomrule
\endlastfoot
Dunlosky, Rawson, Marsh, Nathan \& Willingham (2013) & \textit{Perspectives on Psychological Science in the Public Interest} 14(1), 4--58；DOI: 10.1177/1529100612453266 & 十种学习技术的效用评级：检索练习与分散练习为高效用 \\
Roediger \& Karpicke (2006) & \textit{Psychological Science} 17(3), 249--255；DOI: 10.1111/j.1467-9280.2006.01693.x & 检索练习优于重复阅读（动作一） \\
Karpicke \& Blunt (2011) & \textit{Science} 331(6018), 772--775；DOI: 10.1126/science.1199327 & 检索练习优于概念图；强调"提取"本身 \\
Cepeda, Pashler, Vul, Wixted \& Rohrer (2006) & \textit{Psychological Bulletin} 132(3), 354--380；DOI: 10.1037/0033-2909.132.3.354 & 分散练习的元分析：间隔优于集中（动作二） \\
Rohrer \& Taylor (2007) & \textit{Instructional Science} 35(6), 481--498；DOI: 10.1007/s11251-007-9015-8 & 交错练习优于分块练习（动作三） \\
Chi, Bassok, Lewis, Reimann \& Glaser (1989) & \textit{Cognitive Science} 13(2), 145--182；DOI: 10.1207/s15516709cog1302\_1 & 自我解释与例题学习中的"好学生"行为（动作四） \\
Sweller \& Cooper (1985) & \textit{Cognition and Instruction} 2(1), 59--89；DOI: 10.1207/s1532690xci0201\_3 & 例题效应：先学例题再独立解题（动作五） \\
Ericsson, Krampe \& Tesch-R\"omer (1993) & \textit{Psychological Review} 100(3), 363--406；DOI: 10.1037/0033-295x.100.3.363 & 刻意练习三要素：明确目标、即时反馈、超出舒适区 \\
Bjork (1994) & 载 \textit{Metacognition: Knowing about Knowing}, MIT Press, 185--205 & "合意困难"：让练习变难反而学得更牢 \\
Polya (1945) & \textit{How to Solve It}, Princeton University Press & 解题四步法与本附录"元认知三问" \\
\end{longtable}
}
\normalsize

\textbf{一句总结：}把"我读过了"改成"我\textbf{不看任何提示}能写出核心方程、说出适用条件、给出一个反例"，自学的效率会发生数量级的变化。


\appchapter{附录 J：经典教材逐章对照与延伸阅读}
%======================================================================
% 附录J：经典教材逐章对照与延伸阅读
%======================================================================
% 章标题 \appchapter{附录J：...} 由主文件发出，本文件只提供正文。

附录 E 回答“该买哪本书”，本附录回答一个更具体的问题。\textbf{本书某一章卡住了，去那本书的哪一部分看？}

下面把本书每一章对应到 Chen 与 Fitzpatrick 的章，以及更专门的专著里由谁讲得最清楚。

\begin{warning}{关于章号的使用方式}{}
不同版次的章号会变。本附录的章号取自本节末列出的具体版次（Chen 第 3 版；Fitzpatrick 第 2 版），并已在核对日期当日逐条比对目录。若你手上是别的版次，请按\textbf{主题名}而不是章号查找；本附录每一条都同时给了主题名。
\end{warning}


\section{与附录 E 的分工}

\begin{itemize}
\item \textbf{附录 E}：选书。按水平分级，每条给出“解决本书哪个疑问”。
\item \textbf{本附录}：定位。按下表逐章对照，给出“这一章在哪本书里讲得更透、以及它们讲法上的差别”。
\item \textbf{附录 K}：练手。给出八道跨教材经典例题，并注明出处。
\end{itemize}


\section{七本书的“性格”速写}

\begin{enumerate}
\item \textbf{Chen，\textit{Introduction to Plasma Physics and Controlled Fusion}, 3rd ed.}\\
语言最通俗，物理图像优先，每章有大量数值估计题。适合“先建立直觉”。它的第 9、10 章（Special Plasmas、Plasma Applications）是第 3 版新增内容。与本书第 10--12 章互补。

\item \textbf{Fitzpatrick，\textit{Plasma Physics: An Introduction}, 2nd ed.}\\
推导最完整，且\textbf{全书 HTML 免费在线}。它的第 4 章把流体闭合（Chapman--Enskog、Braginskii）讲得比本书更形式化。第 9 章专门讲磁重联，包括撕裂模的渐近匹配与 Sweet--Parker 模型；本书受篇幅限制只给了量级估计。适合“要把推导补齐”。

\item \textbf{Goldston \& Rutherford，\textit{Introduction to Plasma Physics} (IOP, 1995)。}\\
PPPL 的研究生教材，数值例子贴近聚变实验装置。定位按\textbf{主题}给出（单粒子运动、等离子体流体、碰撞过程、波、动理学），不依赖具体章号。

\item \textbf{Wesson，\textit{Tokamaks} (Oxford University Press)。}\\
托卡马克物理的百科全书式专著。当你想知道“书上这个不稳定性在真实装置上的增长率与阈值是多少”，查它。

\item \textbf{Krall \& Trivelpiece，\textit{Principles of Plasma Physics} (McGraw-Hill, 1973)。}\\
经典“大书”，动理学与波的数学处理最完整，适合当作长期案头参考。

\item \textbf{Hazeltine \& Meiss，\textit{Plasma Confinement} (Addison-Wesley, 1992)。}\\
从约束与输运的角度组织内容，对回旋动理学与输运理论的直觉解释出色。

\item \textbf{Freidberg，\textit{Ideal MHD} (Cambridge University Press, 2014)。}\\
理想 MHD 平衡与稳定性的现代标准参考，能量原理（$\delta W$）讲得最系统；配合本书第 5、9 章使用。
\end{enumerate}


\section{逐章对照表}

\small
\begin{longtable}{@{}p{2.7cm}p{2.5cm}p{2.5cm}p{3.2cm}p{3.0cm}@{}}
\caption{本书各章与经典教材的对照（Fitz 指 Fitzpatrick 第 2 版，Chen 指 Chen 第 3 版）}\label{tab:app-j-map}\\
\toprule
本书 & 主题 & Chen & Fitz & 其他专著 \\
\midrule
\endfirsthead
\toprule
本书 & 主题 & Chen & Fitz & 其他专著 \\
\midrule
\endhead
\midrule
\multicolumn{5}{r}{\small（续下页）}\\
\endfoot
\bottomrule
\endlastfoot
第 1 章 & 数学与统计力学预备知识 & 第 1 章、附录 A（单位、矢量关系） & 无专章，工具散见各章 & Krall \& Trivelpiece 的数学附录 \\
第 2 章 & 什么是等离子体 & 第 1 章 & 第 1 章（判据与参数） & --- \\
第 3 章 & 单粒子运动 & 第 2 章 & 第 2 章（平均化、绝热不变量） & Hazeltine \& Meiss（回旋运动） \\
第 4 章 & 流体描述基础 & 第 3 章 & 第 4 章（Chapman--Enskog、Braginskii） & Bittencourt（输运推导） \\
第 5 章 & 磁流体力学 & 第 3 章、第 6 章 & 第 8 章（磁压、冻结、MHD 波） & Freidberg \textit{Ideal MHD} \\
第 6 章 & 等离子体中的波 & 第 4 章 & 第 5、6、7 章（冷、非均匀、温等离子体） & Stix \textit{Waves in Plasmas} \\
第 7 章 & 动理学理论 & 第 7 章 & 第 3 章、第 7 章（Landau 阻尼、速度空间不稳定性） & Krall \& Trivelpiece \\
第 8 章 & 碰撞与输运 & 第 5 章 & 第 3 章、第 4 章（库仑对数、Rosenbluth 势） & Helander \& Sigmar（新经典输运） \\
第 9 章 & 等离子体不稳定性 & 第 6 章、第 8 章 & 第 7 章、第 9 章（撕裂模、重联） & Freidberg；Wesson \textit{Tokamaks} \\
第 10 章 & 约束与受控核聚变 & 第 9、10 章 & 第 8 章（太阳风、激波） & Wesson \textit{Tokamaks}；Goldston \& Rutherford \\
第 11 章 & 低温工业等离子体 & 第 10 章 & 第 4 章末（Langmuir 鞘层与探针） & Lieberman \& Lichtenberg \\
第 12 章 & 等离子体基础诊断 & 第 10 章 & 第 4 章（探针理论） & Hutchinson \textit{Principles of Plasma Diagnostics} \\
\end{longtable}
\normalsize


\section{每章“外文教材独有视角”}

下面每条是\textbf{本书有意没展开、但在上述教材中讲得更透}的点，也是“读第二本书”时最该先看的地方。

\begin{enumerate}
\item \textbf{第 1 章}：Chen 把单位制与矢量恒等式集中放进附录 A，Fitzpatrick 则把量纲分析直接融入参数定义。两种做法都值得学：前者便于速查，后者让量纲成为物理直觉的一部分。
\item \textbf{第 2 章}：Fitzpatrick 第 1 章用“德拜屏蔽 $\to$ 等离子体参数 $\to$ 碰撞 $\to$ 磁化”的顺序定义等离子体，并配太阳风、电离层的实测数字。比纯理论定义更直观。
\item \textbf{第 3 章}：Fitzpatrick 第 2 章引入“平均化方法”与 Poincaré 不变量的讨论，解释了为什么绝热不变量在长时间尺度上守恒。本书只给结论。
\item \textbf{第 4 章}：Fitzpatrick 第 4 章系统给出 Chapman--Enskog 展开与 Braginskii 输运系数，是本书第 4、8 章之间缺失的那一环。
\item \textbf{第 5 章}：Fitzpatrick 第 8 章用太阳风与 Parker 螺旋、Cowling 反发电机定理把 MHD 放到天体物理语境中。动机比纯聚变视角更宽。
\item \textbf{第 6 章}：Fitzpatrick 把波分成三章（冷、非均匀、温等离子体）。其中非均匀等离子体的 WKB 处理与共振层、射线追踪，是本书完全没有覆盖的内容。
\item \textbf{第 7 章}：Fitzpatrick 第 7 章给出等离子体色散函数 $Z(\zeta)$ 的完整性质表与速度空间不稳定性的统一处理。本书只讨论了 Landau 阻尼的主线。
\item \textbf{第 8 章}：Fitzpatrick 第 3 章的 Landau 碰撞算符与 Rosenbluth 势比本书形式化得多。Chen 第 5 章则给出与实验对照的扩散、电阻率讨论。一个补数学，一个补实验。
\item \textbf{第 9 章}：Fitzpatrick 第 9 章用渐近匹配（asymptotic matching）求解撕裂模，并讨论 Sweet--Parker 与 plasmoid 不稳定性。这是理解现代重联研究的标准入口。
\item \textbf{第 10 章}：Goldston \& Rutherford 与 Wesson 提供真实装置尺度的数值（密度、温度、约束时间、阈值），可以把本书的劳森判据估算校准到工程现实。
\item \textbf{第 11 章}：Fitzpatrick 第 4 章末把鞘层与探针作为流体理论的直接应用；Lieberman \& Lichtenberg 则从刻蚀、沉积等工艺角度组织内容。这是“低温等离子体工程”的标准参考。
\item \textbf{第 12 章}：Hutchinson 按“测什么物理量”（密度、温度、磁场、电流分布）组织诊断。与本书按技术（探针、光谱、干涉、汤姆孙散射）组织的方式互为补充，适合做诊断选型的依据。
\end{enumerate}


\section{进阶专著速查}

{\small
\begin{longtable}{@{}p{4.4cm}p{5.2cm}p{4.9cm}@{}}
\caption{进入专门方向时的参考书}\label{tab:app-j-advanced}\\
\toprule
专著 & 出版信息 & 什么时候读它 \\
\midrule
\endfirsthead
\toprule
专著 & 出版信息 & 什么时候读它 \\
\midrule
\endhead
\midrule
\multicolumn{3}{r}{\small（续下页）}\\
\endfoot
\bottomrule
\endlastfoot
Wesson, \textit{Tokamaks} & Oxford University Press（初版 1987） & 处理托卡马克的具体不稳定性、约束与运行极限时 \\
Krall \& Trivelpiece, \textit{Principles of Plasma Physics} & McGraw-Hill, 1973 & 需要最完整的动理学与波的数学处理时 \\
Hazeltine \& Meiss, \textit{Plasma Confinement} & Addison-Wesley, 1992 & 研究约束、回旋动理学与输运理论时 \\
Freidberg, \textit{Ideal MHD} & Cambridge University Press, 2014 & 系统学习 MHD 平衡与稳定性、能量原理时 \\
Goedbloed \& Poedts, \textit{Principles of Magnetohydrodynamics} & Cambridge University Press, 2004 & 需要 MHD 谱理论（连续谱、阿尔芬本征模）时 \\
Bittencourt, \textit{Fundamentals of Plasma Physics} & Pergamon, 1986；Springer 第 3 版, 2004 & 需要比本书更详细的输运与波推导时 \\
Hutchinson, \textit{Principles of Plasma Diagnostics} & Cambridge University Press（初版 1987；第 2 版 2002） & 设计或解释诊断测量时 \\
Helander \& Sigmar, \textit{Collisional Transport in Magnetized Plasmas} & Cambridge University Press, 2002 & 研究新经典输运与杂质输运时 \\
Stix, \textit{Waves in Plasmas} & American Institute of Physics, 1992 & 需要波支分类与射频加热的系统处理时 \\
Bellan, \textit{Fundamentals of Plasma Physics} & Cambridge University Press, 2006 & 需要从第一性原理严格推导基础方程时 \\
\end{longtable}
}


\section{免费在线资源}

\begin{itemize}
\item \textbf{Fitzpatrick 教材书页}：\url{https://farside.ph.utexas.edu/Books/Plasma2/Plasma2.html}\\
作者提供的第 2 版在线版本（HTML），与纸质书内容对应，适合对照本附录中的章号。
\item \textbf{Fitzpatrick 等离子体物理讲义}\\
\url{https://farside.ph.utexas.edu/teaching/plasma/Plasmahtml/Plasmahtml.html}\\
单学期课程的讲义，篇幅更短，适合作为读本书时的并行材料。
\item \textbf{粒子模拟示例程序}：\url{https://farside.ph.utexas.edu/teaching/plasma/pic.tgz}\\
作者配套的 PIC 程序（含双流不稳定性算例），可用于验证本书第 9 章的线性增长率。
\end{itemize}

上述三个链接均为作者个人主页公开资源，核对日期见下节。


\section{元数据与核对说明}

本附录的章号与出版信息按以下方式核对（核对日期：2026-10-11）：

\begin{itemize}
\item Chen 第 3 版目录：经图书馆编目与出版方书目信息交叉核对，确认第 3 版\textbf{没有}单独的“碰撞效应”章。碰撞内容散见于第 5 章（扩散与电阻率）与第 7 章（动理学理论）。
\item Fitzpatrick 第 2 版目录：取自作者官方书页（本附录“免费在线资源”第一条），逐章逐节核对，共 9 章。
\item Goldston \& Rutherford：只核实到“六大部分”的主题划分，因此本附录对其采用\textbf{主题级}定位，不给章号。
\item 进阶专著的出版社、版次与年份：经开放图书馆（OpenLibrary）书目数据核对。
\item 本附录不引用任何未获授权的电子副本；Fitzpatrick 的在线版本为作者本人公开提供。
\end{itemize}

\begin{tip}{对照表用法}{tip-appj-howto}
\textbf{怎么看这张对照表：}不要按“章号对齐”去读书，而要按\textbf{问题}去读。当你卡在某个推导时，先在表中找到本书章节所在的行，再看哪一列的主题名与你的问题最接近，然后去那一章找对应的主题名。
\end{tip}


\appchapter{附录 K：跨教材经典例题精选}
%======================================================================
% 附录K：跨教材经典例题精选（注明出处）
%======================================================================
% 章标题 \appchapter{附录K：...} 由主文件发出，本文件只提供正文。

本附录从 Chen、Fitzpatrick、Goldston \& Rutherford 等经典教材中挑八道\textbf{最能检验基础图像}的题，逐题给出完整解答与出处。

选题只看三条：只用本书第 1--12 章的工具，答案能\textbf{独立复算}，而且背后都有一个「判断型」的物理结论，不只是代数练习。

\textbf{使用建议：}每道题先合书独立做一遍（动作五），再看解答。八道题覆盖八个不同的物理环节，适合按附录 I 的「交错练习」方式打乱使用。

\begin{tip}{复算与自查}{tip-appk-recalc}
若你的结果与文中差一个量级，几乎总是以下几个原因之一：用错了温度（电子 / 离子、eV / K）、漏掉平方根、或把 SI 与 CGS 混用。建议对照附录 F 的高频易错点清单一起看。
\end{tip}


\section{德拜屏蔽与等离子体参数}

\begin{example}{德拜长度、$N_D$ 与屏蔽势的衰减}{}
\textbf{题目：}某氢等离子体 $T_e=10\ \mathrm{eV}$、$n=10^{18}\ \mathrm{m^{-3}}$。
（a）求德拜长度 $\lambda_D$。
（b）求德拜球内的粒子数 $N_D$，并判断能否用流体（连续介质）图像描述。
（c）把测试电荷放在原点，$r=5\lambda_D$ 处的屏蔽势是裸库仑势的百分之几？

\textbf{解：}
（a）$\lambda_D=\sqrt{\dfrac{\eps_0 k_BT_e}{ne^2}}=7430\sqrt{\dfrac{T_e[\mathrm{eV}]}{n[\mathrm{m^{-3}}]}}\ \mathrm{m}
=7430\sqrt{\dfrac{10}{10^{18}}}=2.35\times10^{-5}\ \mathrm{m}\approx23.5\ \mathrm{\mu m}$。

（b）$N_D=\dfrac{4\pi}{3}n\lambda_D^3
=\dfrac{4\pi}{3}\times10^{18}\times(2.3508\times10^{-5})^3\approx5.4\times10^{4}$。
$N_D\gg1$，德拜球内粒子数极多，统计描述成立，流体图像在这里是合法的。

（c）屏蔽势 $\phi(r)=\dfrac{q}{4\pi\eps_0 r}\ee^{-r/\lambda_D}$，比值即 $\ee^{-5}=6.7\times10^{-3}$，只有 $0.67\%$。

{\footnotesize
\textbf{出处：}Chen《Introduction to Plasma Physics and Controlled Fusion》第 3 版第 1 章（德拜屏蔽）。
Fitzpatrick 第 2 版第 1 章 Debye Shielding。
}

\textbf{常见错误：}用离子温度算 $\lambda_D$（判据由电子决定）；漏掉 $N_D$ 中的 $4\pi/3$ 因子；把「$\lambda_D$ 内只有一个粒子」当成判据（正确判据是 $N_D\gg1$）。
\end{example}


\section{磁镜与损失锥}

\begin{example}{磁镜比 $R=10$ 的临界角与损失锥比例}{}
\textbf{题目：}磁镜的磁镜比 $R=B_{\max}/B_{\min}=10$。
（a）求临界投掷角 $\theta_c$。
（b）各向同性速度分布中，有多少比例的粒子落在损失锥内？
（c）若把 $R$ 提高到 $100$，该比例变为多少？

\textbf{解：}
（a）由 $\sin^2\theta_c=B_{\min}/B_{\max}=1/R$ 得。
$\theta_c=\arcsin\sqrt{1/10}=18.43^\circ$。

（b）各向同性分布中投掷角概率密度正比于 $\sin\theta$，故损失锥内的比例为。
$P=\displaystyle\int_0^{\theta_c}\sin\theta\,\diff\theta=1-\cos\theta_c
=1-\sqrt{1-\tfrac1R}=1-\sqrt{0.9}=5.13\%$。

（c）$R=100$ 时 $P=1-\sqrt{0.99}=0.50\%$：磁镜比提高十倍，损失锥比例降到约十分之一。

{\footnotesize
\textbf{出处：}Chen 第 2 章（单粒子运动，磁镜）。
Fitzpatrick 第 2 版第 2 章 Magnetic Mirrors。
}

\textbf{常见错误：}写成 $\sin\theta_c=1/R$（漏掉开方）；用 $v_\parallel/v_\perp$ 而不是 $\sin^2\theta$ 表达判据；把损失锥比例当成 $\theta_c/90^\circ$ 的线性比例。
\end{example}


\section{梯度漂移与曲率漂移的方向}

\begin{example}{$\grad B$ 漂移的方向判别与电荷分离}{}
\textbf{题目：}设磁场沿 $z$ 方向、大小沿 $x$ 线性减弱：$\vect{B}=B_0(1-x/L)\unitvec{z}$，$B_0>0$，$L>0$。
（a）求 $\grad B$，并判断离子与电子的 $\grad B$ 漂移方向。
（b）说明在托卡马克中，$\grad B$ 漂移与曲率漂移为什么会导致电荷分离，以及它是如何被抑制的。

\textbf{解：}
（a）$\grad B=-\dfrac{B_0}{L}\unitvec{x}$，于是
$\vect{B}\times\grad B=B_0\left(1-\dfrac{x}{L}\right)\unitvec{z}\times\left(-\dfrac{B_0}{L}\unitvec{x}\right)
=-\dfrac{B_0^2}{L}\left(1-\dfrac{x}{L}\right)\unitvec{y}$。
代入 $\vect{v}_{\grad B}=\dfrac{mv_\perp^2}{2qB}\dfrac{\vect{B}\times\grad B}{B^2}$：
离子（$q>0$）沿 $-\unitvec{y}$ 漂移，电子（$q<0$）沿 $+\unitvec{y}$ 漂移。方向相反，与电荷符号有关。

（b）在托卡马克中 $B\propto1/R$，$\grad B$ 指向内侧。这里最容易错的一步是\textbf{把曲率矢量与漂移方向混为一谈}：曲率矢量 $\kappa=(\unitvec{b}\cdot\grad)\unitvec{b}$ 确实指向曲率中心（环向场下 $\kappa=-\unitvec{R}/R_c$），但曲率漂移沿 $\unitvec{b}\times\kappa$，\emph{垂直于}曲率矢量：
\[
\vect{v}_R=\frac{mv_\parallel^2}{qB}\,(\unitvec{b}\times\kappa),
\qquad
\unitvec{b}\times\kappa=\unitvec{\varphi}\times\left(-\frac{\unitvec{R}}{R_c}\right)=+\frac{\unitvec{z}}{R_c}.
\]
漂移方向是竖直的（正离子向上、电子向下），\textbf{不是}指向曲率中心。
它与 $\grad B$ 漂移的\textbf{竖直分量同向}（后者也沿 $\pm\unitvec{z}$），且都因含 $1/q$ 而对电子、离子反向，于是等离子体上下两侧积累异号电荷，形成竖直电场；该电场与环向磁场产生 $\vect{E}\times\vect{B}$ 漂移，把粒子向外推。
抑制机制是\textbf{旋转变换}（磁力线的环向扭曲）：粒子沿磁力线运动时上下交替，短接了电荷分离。托卡马克因此需要足够大的等离子体电流，「磁通管」图像中才会出现 Pfirsch--Schl\"uter 电流。

{\footnotesize
\textbf{出处：}Chen 第 2 章（单粒子运动，非均匀磁场中的漂移）。
Fitzpatrick 第 2 版第 2 章 Magnetic Drifts。

托卡马克部分见 Wesson《Tokamaks》（按主题）。
}

\textbf{常见错误：}认为 $\grad B$ 漂移与电荷符号无关。$\vect{B}\times\grad B$ 确实与 $q$ 无关，但漂移速度中含 $1/q$，所以方向相反。另一种错法是把曲率漂移的方向当作曲率矢量的方向：漂移 $\propto\unitvec{b}\times\kappa$，与 $\kappa$ 垂直，两者不可互换。还有人把曲率漂移与 $\grad B$ 漂移在本例的几何下判成反向，其实它们同向；也有人以为电荷分离无法消除。
\end{example}


\section{Alfvén 速度与磁冻结}

\begin{example}{Alfvén 速度的量级与磁场扰动的传播时间}{}
\textbf{题目：}氘等离子体 $n=10^{20}\ \mathrm{m^{-3}}$、$B=5\ \mathrm{T}$。
（a）求 Alfvén 速度 $v_A$。（b）它是光速的百分之几？
（c）磁场扰动沿 $1\ \mathrm{m}$ 传播需要多久？（d）若密度增大到 $10^{22}\ \mathrm{m^{-3}}$，$v_A$ 变为多少？

\textbf{解：}
（a）质量密度由离子提供：$\rho\approx nm_D=10^{20}\times3.368\times10^{-27}=3.37\times10^{-7}\ \mathrm{kg\,m^{-3}}$。
$v_A=\dfrac{B}{\sqrt{\mu_0\rho}}
=\dfrac{5}{\sqrt{1.2566\times10^{-6}\times3.368\times10^{-7}}}
=7.69\times10^{6}\ \mathrm{m\,s^{-1}}$。

（b）$v_A/c=7.69\times10^{6}/2.998\times10^{8}=2.6\%$。这里只能得出一个结论：\textbf{非相对论 Alfvén 速度公式适用}（位移电流可以忽略），\textbf{不能据此断定磁冻结成立}。磁冻结要求理想欧姆定律 $\vect{E}+\vect{v}\times\vect{B}=0$ 成立，至少要磁雷诺数 $R_m=\mu_0Lv/\eta\gg1$，还要检查霍尔项与电子压强项是否可丢（第 5 章「MHD 适用条件：三道独立的门槛」）。本题没有给温度、电阻率与流动尺度，用现有输入无法判断：\emph{同一个} $v_A/c$ 在强碰撞的电阻性等离子体里成立，在理想 MHD 等离子体里也成立。作为对照：取 $L=1\ \mathrm{m}$、$\eta(10\ \mathrm{eV})=1.64\times10^{-5}\ \Omega\cdot\mathrm{m}$（K.6 已算），则
\[
R_m=\frac{\mu_0Lv_A}{\eta}
=\frac{1.2566\times10^{-6}\times1\times7.69\times10^{6}}{1.64\times10^{-5}}
\approx5.9\times10^{5}\gg1,
\]
此时才能谈磁冻结。而电阻扩散时标 $\tau_\eta=\mu_0L^2/\eta\approx0.076\ \mathrm{s}$ 给出了一个定量边界：在 $L\sim1\ \mathrm{m}$ 的尺度上，静止导电流体的磁场正是以这个时标扩散掉的（见 K.6）。

（c）$t=L/v_A=1/7.69\times10^{6}=1.30\times10^{-7}\ \mathrm{s}$。

（d）$v_A\propto1/\sqrt{\rho}\propto1/\sqrt{n}$：密度增大 100 倍，$v_A$ 减小到 $1/10$，即 $7.7\times10^{5}\ \mathrm{m\,s^{-1}}$。

{\footnotesize
\textbf{出处：}Fitzpatrick 第 2 版第 8 章（MHD 流体：磁压、磁通冻结、MHD 波）。
Chen 第 4 章（等离子体中的波）。
}

\textbf{常见错误：}用电子质量计算 $\rho$，而惯性由离子提供。忘记 $\mu_0$ 的平方根。把 $v_A$ 与离子声速 $c_s$ 混用：后者由电子温度决定，与 $B$ 无关。把 $v_A/c\ll1$ 当成磁冻结的充分条件，而它只保证非相对论公式适用；$R_m=\mu_0Lv/\eta\gg1$ 才是丢掉电阻项的条件。
\end{example}


\section{Landau 阻尼率}

\begin{example}{朗缪尔波的 Landau 阻尼率与每周期衰减}{}
\textbf{题目：}朗缪尔波满足 $k\lambda_D=0.3$。
（a）用 $\dfrac{\gamma}{\omega}=-\sqrt{\dfrac{\pi}{8}}\,(k\lambda_D)^{-3}
\exp\!\left(-\dfrac{1}{2k^2\lambda_D^2}-\dfrac{3}{2}\right)$ 求阻尼率。
（b）估算一个波周期内振幅衰减的百分比。
（c）为什么 $k\lambda_D\ll1$ 时阻尼可以忽略？

\textbf{解：}
（a）$(k\lambda_D)^{-3}=0.3^{-3}=37.04$，
$\exp[-1/(2\times0.09)]=\ee^{-5.556}=3.87\times10^{-3}$，$\ee^{-3/2}=0.2231$，
故 $\dfrac{\gamma}{\omega}=-0.6267\times37.04\times3.87\times10^{-3}\times0.2231=-2.00\times10^{-2}$。

（b）一个周期内振幅因子为 $\exp\!\left(2\pi\dfrac{\gamma}{\omega}\right)=\ee^{-0.126}=0.882$，即衰减约 $11.8\%$。

（c）指数因子 $\exp[-1/(2k^2\lambda_D^2)]$ 随 $k\lambda_D$ 减小而\textbf{双重指数}衰减：$k\lambda_D=0.1$ 时它已是 $\ee^{-50}\approx2\times10^{-22}$。物理上，长波相速度远大于热速度，能与波共振的粒子数指数稀少，因此无碰撞阻尼可以忽略。

{\footnotesize
\textbf{出处：}Fitzpatrick 第 2 版第 7 章 Landau Damping。
Chen 第 7 章（动理学理论）。
}

\textbf{常见错误：}把 $\gamma/\omega$ 直接当成「每周期衰减比例」（应乘以 $2\pi$）；把 $\ee^{-3/2}$ 漏掉；误认为 Landau 阻尼来自碰撞；它是无碰撞的相混合，不增加熵（在线性阶段）。
\end{example}


\section{Spitzer 电阻率与磁场扩散时间}

\begin{example}{从温度估算电阻率与磁场扩散时间}{}
\textbf{题目：}某托卡马克等离子体 $T=1\ \mathrm{keV}$、$Z=1$、$\ln\Lambda=17$，特征长度 $L=1\ \mathrm{m}$。
（a）用 $\eta=5.2\times10^{-5}\dfrac{Z\ln\Lambda}{T[\mathrm{eV}]^{3/2}}\ \Omega\cdot\mathrm{m}$ 求电阻率。
（b）求磁场扩散时间 $\tau_\eta=\mu_0L^2/\eta$。
（c）若温度降到 $10\ \mathrm{eV}$（$\ln\Lambda=10$），$\eta$ 与 $\tau_\eta$ 各变为多少？

\textbf{解：}
（a）$T^{3/2}=1000^{1.5}=3.16\times10^{4}$，
$\eta=5.2\times10^{-5}\times17/3.16\times10^{4}=2.80\times10^{-8}\ \Omega\cdot\mathrm{m}$。

（b）$\tau_\eta=\dfrac{1.2566\times10^{-6}\times1}{2.80\times10^{-8}}=44.9\ \mathrm{s}$。

（c）$10^{1.5}=31.6$，$\eta=5.2\times10^{-5}\times10/31.6=1.64\times10^{-5}\ \Omega\cdot\mathrm{m}$（约为原来的 $590$ 倍），
$\tau_\eta=1.2566\times10^{-6}/1.64\times10^{-5}=7.6\times10^{-2}\ \mathrm{s}$。

\textbf{结论：}温度提高两个数量级，$\eta$ 下降约 $590$ 倍，磁场扩散时间相应延长近三个数量级。所以\textbf{在给定尺度与流速下}，磁雷诺数 $R_m=\mu_0Lv/\eta$ 随温度升高而增大：高温聚变等离子体是磁冻结的好近似，低温器件中磁场可以迅速扩散。（注意 $R_m$ 还含 $L$ 与 $v$，扩散时标本身只由 $\eta$ 与尺度决定。）

{\footnotesize
\textbf{出处：}Chen 第 5 章（扩散与电阻率）。
Fitzpatrick 第 2 版第 3、4 章（碰撞与输运）。
}

\textbf{常见错误：}把开尔文直接代入公式（应为 eV）；漏掉 $T^{3/2}$；忽略 $\ln\Lambda$ 的温度依赖（它变化很慢，但并非常数）。
\end{example}


\section{玻姆判据与鞘层厚度}

\begin{example}{离子声速、鞘层厚度与壁面离子通量}{}
\textbf{题目：}氩等离子体 $T_e=3\ \mathrm{eV}$、$n_0=10^{16}\ \mathrm{m^{-3}}$，器壁偏压 $V_w=100\ \mathrm{V}$。
（a）求玻姆速度（离子声速）$c_s$。（b）求\emph{体}等离子体的德拜长度。
（c）用 $s\approx\dfrac{\sqrt2}{3}\lambda_D(n_s)\left(\dfrac{2V_w}{T_e}\right)^{3/4}$ 估算 Child--Langmuir 鞘层厚度（$\lambda_D(n_s)$ 为鞘边密度处的德拜长度）。
（d）求壁面离子通量，并说明它与 (c) 的厚度是否守同一个电流边界条件。

\textbf{解：}
（a）$c_s=\sqrt{\dfrac{k_BT_e}{M_{\mathrm{Ar}}}}
=\sqrt{\dfrac{3\times1.602\times10^{-19}}{6.68\times10^{-26}}}=2.68\times10^{3}\ \mathrm{m\,s^{-1}}$。

（b）\emph{体}等离子体的德拜长度 $\lambda_D(n_0)=7430\sqrt{3/10^{16}}=1.29\times10^{-4}\ \mathrm{m}=0.129\ \mathrm{mm}$。
注意鞘层计算要用\textbf{鞘边}密度 $n_s$ 对应的德拜长度，见 (c)。

（c）鞘层从鞘边密度 $n_s=0.61n_0$ 起步，所以厚度估算式里的 $\lambda_D$ 必须取 $n_s$ 处的值。
$\lambda_D(n_s)=\lambda_D(n_0)/\sqrt{0.61}=1.29\times10^{-4}/0.781=1.65\times10^{-4}\ \mathrm{m}$。
$(2\times100/3)^{0.75}=66.7^{0.75}=23.3$，
$s=0.471\times1.65\times10^{-4}\times23.3=1.81\times10^{-3}\ \mathrm{m}\approx1.8\ \mathrm{mm}$。
注意 $s\gg\lambda_D(n_s)$（约 11 倍）：偏压远大于电子温度时，鞘层远厚于德拜长度。
若误用体密度算出的 $\lambda_D(n_0)$，厚度会被低估 $\sqrt{n_0/n_s}=1.28$ 倍（得到 1.4 mm），并使厚度与 (d) 的通量不守同一个电流边界条件。

（d）鞘边密度 $n_s\approx0.61n_0=6.1\times10^{15}\ \mathrm{m^{-3}}$。
$\Gamma=n_sc_s=6.1\times10^{15}\times2.68\times10^{3}=1.6\times10^{19}\ \mathrm{m^{-2}s^{-1}}$，
即壁面离子电流密度 $j=e\Gamma\approx2.62\ \mathrm{A\,m^{-2}}$。用 (c) 的厚度反代平面无碰撞 Child--Langmuir 关系
\[
j=\frac49\eps_0\sqrt{\frac{2e}{M}}\frac{V_w^{3/2}}{s^2}
\]
得到 $j\approx2.62\ \mathrm{A\,m^{-2}}$（由该式反解鞘厚得 $1.813\ \mathrm{mm}$，与估算式相差 0.2\%）。厚度与通量守同一个电流边界条件。
每个离子在鞘层中被加速获得约 $100\ \mathrm{eV}$ 能量，这正是离子轰击在刻蚀与溅射中起作用的原因。

{\footnotesize
\textbf{出处：}Fitzpatrick 第 2 版第 4 章末 Langmuir Sheaths 与 Langmuir Probes。
Lieberman \& Lichtenberg《Principles of Plasma Discharges and Materials Processing》（按主题：鞘层与玻姆判据）。
}

\textbf{常见错误：}用离子温度计算 $c_s$，而玻姆判据由电子温度决定。把鞘层厚度当成 $\lambda_D$。把\emph{体}密度 $n_0$ 的德拜长度用到鞘边，厚度会被低估 $\sqrt{n_0/n_s}=1.28$ 倍。厚度与壁面通量若用不同密度，两者就不守同一个 Child--Langmuir 电流边界条件。忘记 $n_s\approx0.61n_0$ 这个鞘边密度。
\end{example}


\section{Sweet--Parker 重联速率}

\begin{example}{Sweet--Parker 重联为什么「太慢」}{}
\textbf{题目：}日冕电流片：$T=86\ \mathrm{eV}$（约 $10^{6}\ \mathrm{K}$）、$\ln\Lambda=20$、$L=10^{5}\ \mathrm{m}$、$v_A=10^{5}\ \mathrm{m\,s^{-1}}$。
（a）求电阻率 $\eta$ 与 Lundquist 数 $S=\mu_0Lv_A/\eta$。
（b）求入流速度 $v_{\mathrm{in}}=v_A/\sqrt{S}$ 与重联时间 $\tau_{\mathrm{SP}}=L/v_{\mathrm{in}}$。
（c）与太阳耀斑观测时标（$10^{2}$--$10^{3}\ \mathrm{s}$）比较，能得出什么结论？

\textbf{解：}
（a）$86^{1.5}=7.98\times10^{2}$，
$\eta=5.2\times10^{-5}\times20/7.98\times10^{2}=1.30\times10^{-6}\ \Omega\cdot\mathrm{m}$；
$S=\dfrac{1.2566\times10^{-6}\times10^{5}\times10^{5}}{1.30\times10^{-6}}=9.6\times10^{9}$。

（b）$v_{\mathrm{in}}=v_A/\sqrt{S}=10^{5}/9.82\times10^{4}=1.02\ \mathrm{m\,s^{-1}}$，
$\tau_{\mathrm{SP}}=10^{5}/1.02=9.8\times10^{4}\ \mathrm{s}\approx27\ \mathrm{h}$。

（c）$27$ 小时比观测到的耀斑时标慢约 $10^{2}$--$10^{3}$ 倍。这就是著名的「Sweet--Parker 太慢」问题：全局电阻率太低，重联被限制在极薄的电流片内。解决方向是\textbf{局域化}：撕裂模不稳定性把电流片撕裂成许多小尺度结构（plasmoid 不稳定性），使有效重联率大幅提高。

\textbf{补充关系：}磁场扩散时间 $\tau_\eta=\mu_0L^2/\eta=\sqrt{S}\,\tau_{\mathrm{SP}}=9.6\times10^{9}\ \mathrm{s}$（约 300 年），比 $\tau_{\mathrm{SP}}$ 大 $\sqrt{S}\approx10^{5}$ 倍，这就排除了「整体扩散」的解释，也说明重联必须发生在薄电流片中。

{\footnotesize
\textbf{出处：}Fitzpatrick 第 2 版第 9 章（磁重联：Sweet--Parker 模型、撕裂模、plasmoid 不稳定性）。
}

\textbf{常见错误：}把入流速度取成 $v_A$（$v_A$ 是入流在上游的极限值，重联率取决于 $\sqrt{S}$ 的压制）；漏掉 $S$ 定义中的 $L$ 或 $v_A$；把 $\tau_{\mathrm{SP}}$ 与 $\tau_\eta$ 混为一谈。
\end{example}


\section{八道题与本书章节的对应}

\small
\begin{longtable}{@{}p{1.2cm}p{4.5cm}p{4.3cm}p{4.9cm}@{}}
\caption{附录 K 例题与本书章节、经典教材的对应}\label{tab:app-k-map}\\
\toprule
编号 & 主题 & 本书章节 & 主要出处 \\
\midrule
\endfirsthead
\toprule
编号 & 主题 & 本书章节 & 主要出处 \\
\midrule
\endhead
\midrule
\multicolumn{4}{r}{\small（续下页）}\\
\endfoot
\bottomrule
\endlastfoot
K.1 & 德拜屏蔽与 $N_D$ & 第 2 章 & Chen 第 1 章；Fitzpatrick 第 1 章 \\
K.2 & 磁镜与损失锥 & 第 3 章 & Chen 第 2 章；Fitzpatrick 第 2 章 \\
K.3 & 梯度漂移与曲率漂移 & 第 3 章 & Chen 第 2 章；Fitzpatrick 第 2 章 \\
K.4 & Alfvén 速度与磁冻结 & 第 5 章 & Fitzpatrick 第 8 章；Chen 第 4 章 \\
K.5 & Landau 阻尼 & 第 7 章 & Fitzpatrick 第 7 章；Chen 第 7 章 \\
K.6 & Spitzer 电阻率与扩散 & 第 8 章 & Chen 第 5 章；Fitzpatrick 第 3、4 章 \\
K.7 & 玻姆判据与鞘层 & 第 11 章 & Fitzpatrick 第 4 章；Lieberman \& Lichtenberg \\
K.8 & Sweet--Parker 重联 & 第 9 章 & Fitzpatrick 第 9 章 \\
\end{longtable}
\normalsize

\textbf{做完这八道题之后：}回到各章的自学检查清单（\texttt{checklistbox}）逐条自评，并把做错的题目按附录 I 的三栏法记入错题本。若你发现某道题「看答案全懂、合上书不会」，说明缺的是检索练习而不是理解。按附录 I 的动作五，隔一天再做一遍。


\appchapter{附录 L：分章自测题库（含答案与提示）}
%======================================================================
% 附录L：分章自测题库（含答案与提示）
%======================================================================
% 章标题 \appchapter{附录L：...} 由主文件发出，本文件只提供正文。
% 设计约定：每章固定 6 题 = 2 概念 + 2 推导 + 2 估算；每题给「答案要点」。
% 版权约定：只借鉴经典教材的题型与难度分级，题目、数值由本书自拟并独立复算。

\newcommand{\tkans}[1]{\par\vspace{0.5pt}{\small\textbf{答案：}#1}\par}
\newcommand{\tksrc}[1]{{\small\textbf{取材：}#1}\par\vspace{1pt}}

本附录给出一份\textbf{可执行的分章自测题库}：第 1--12 章每章 6 题，固定为 \textbf{2 道概念题 + 2 道推导题 + 2 道估算题}，每题附「答案要点」（给思路与关键数字，而不是只给一个数）。

它与另外两个附录的分工是：附录 K 是\textbf{范例}，给八道跨教材经典题与完整解答，用来学「怎么写」。附录 L 是\textbf{自测}，先合书独立做再看要点，用来查「会不会」。附录 I 是\textbf{策略}，管怎么安排复习、怎么处理错题。

\tksrc{每节的题型与难度分级参照 Chen 与 Fitzpatrick 教材对应章的习题类型（逐章对照见附录 J 的表 \ref{tab:app-j-map} 与本节末的来源表）；\textbf{所有题目与数值均由本书自拟，并已用脚本独立复算}，不复制任何原书题目。}

\begin{tip}{题库用法}{tip-appl-howto}
用法建议：先按附录 I 的「合书复述」独立做完一章的 6 题，再对答案要点。若一题的结果与文中\textbf{差一个量级}，先别怀疑公式，按附录 F 的顺序排查：温度用错（电子 / 离子、eV / K）、漏掉平方根、SI 与 CGS 混用、把 $kT$ 与 $T$ 等同。数值题一律给到两位有效数字即可，重点是量级与结论方向。
\end{tip}


\section{第 1 章：数学与统计力学预备知识}
\tksrc{Chen 第 1 章与附录 A（单位、矢量关系）、Fitzpatrick 各章散见的数学工具。}

\begin{enumerate}
\item （概念）德拜屏蔽与高斯定理矛盾吗？为什么屏蔽后的势不再是库仑势？
\tkans{不矛盾。高斯定理对\textbf{总场}成立：测试电荷周围的反号屏蔽云贡献的通量与点电荷通量部分抵消，$r\gg\lambda_D$ 时净通量趋于零。库仑势只是无屏蔽极限（$N_D\to\infty$）；德拜势 $(\frac{q}{4\pi\eps_0 r})e^{-r/\lambda_D}$ 是把屏蔽云的平均场计入后的结果。}

\item （概念）写出麦克斯韦速率分布的三个前提，并说明为什么最概然速率、平均速率、均方根速率三者不相等。
\tkans{前提：无外场、热平衡、非简并（经典统计）。速率分布 $f(v)\propto v^2e^{-mv^2/2k_BT}$ 不是关于最概然值的对称分布，$v^2$ 因子把分布向高速侧拉长，因此 $v_{th}:v_{avg}:v_{rms}=1:1.128:1.225$（三者依次增大）。}

\item （推导）由麦克斯韦分布导出 $v_{th}=\sqrt{2k_BT/m}$ 与 $v_{avg}=\sqrt{8k_BT/\pi m}$，并计算 $T_e=10\ \mathrm{eV}$ 时电子的两个数值。
\tkans{$v_{th}$ 由 $f(v)$ 的极大值给出，$v_{avg}$ 由分布的一阶矩给出，二者不相等。电子：$v_{th}=\sqrt{2T_e/m_e}=1.9\times10^{6}\ \mathrm{m/s}$，$v_{avg}=(2/\sqrt{\pi})\,v_{th}=2.1\times10^{6}\ \mathrm{m/s}$。注意两者相差 13\%，口算时不要混用。}

\item （推导）证明理想气体 $\langle\frac12mv^2\rangle=\frac32k_BT$，并说明「$1\ \mathrm{eV}$ 对应多少开尔文」。
\tkans{三个平动自由度各贡献 $\frac12k_BT$（能量均分）。由 $k_BT=e\cdot1\ \mathrm{V}$ 得 $1\ \mathrm{eV}=e/k_B=1.602\times10^{-19}/1.381\times10^{-23}=11605\ \mathrm{K}$。聚变等离子体「$10\ \mathrm{keV}$」即 $1.16\times10^{8}\ \mathrm{K}$，温度极高但粒子数密度低，压强仍在工程可及范围。}

\item （估算）$1\ \mathrm{Pa}$、$300\ \mathrm{K}$ 的中性气体数密度是多少？与 $n=10^{18}\ \mathrm{m^{-3}}$ 的「等离子体」相比，真空度谁更高？
\tkans{$n=p/k_BT=1/(1.381\times10^{-23}\times300)=2.4\times10^{20}\ \mathrm{m^{-3}}$，比 $10^{18}$ 高 240 倍。也就是说 $10^{18}\ \mathrm{m^{-3}}$ 的等离子体所处的本底真空必须优于 $10^{-2}\ \mathrm{Pa}$ 量级，否则中性粒子碰撞会主导，这正是实验室装置要抽高真空的原因。}

\item （估算）若 $n=10^{15}\ \mathrm{m^{-3}}$、$T=1\ \mathrm{eV}$，求 $\lambda_D$ 与德拜球内粒子数 $N_D$，并判断统计描述是否成立。
\tkans{$\lambda_D=7430\sqrt{T[\mathrm{eV}]/n[\mathrm{m^{-3}}]}=7430\sqrt{10^{-15}}=2.4\times10^{-4}\ \mathrm{m}$（0.24 mm）；$N_D=\frac{4\pi}{3}n\lambda_D^3\approx5.4\times10^{4}\gg1$。统计描述成立，但 $\lambda_D$ 已到毫米量级，装置尺度必须远大于 0.24 mm 才能谈「等离子体」（判据 $\lambda_D\ll L$ 与 $N_D\gg1$ 缺一不可）。}
\end{enumerate}


\section{第 2 章：什么是等离子体}
\tksrc{Chen 第 1 章（判据与参数）、Fitzpatrick 第 1 章。}

\begin{enumerate}
\item （概念）等离子体的三条判据各解决什么问题？
\tkans{$\lambda_D\ll L$：保证准中性、外场能被屏蔽（「集体」的前提）；$N_D\gg1$：保证德拜球内有足够多粒子，统计（连续介质）描述成立；$\omega_{pe}\tau>1$（$\tau$ 为碰撞时间）：保证集体振荡在两次碰撞之间能完成，等离子体行为不被碰撞抹平。三者分别对应空间尺度、统计性、时间尺度。}

\item （概念）为什么弱电离的辉光放电和完全电离的日冕都算等离子体，而火焰和电离层 D 层边缘要谨慎？
\tkans{判据看的是\textbf{集体性}而不是电离度：辉光放电 $n$ 大、$T_e$ 高，$\lambda_D$ 小、$N_D$ 大，仍满足三条判据；日冕虽稀薄，但尺度 $L$ 极大，$\lambda_D\ll L$ 成立。火焰电离度极低、且以化学反应与对流为主，集体行为被淹没；D 层边缘 $n$ 很低而碰撞频繁（$\omega_{pe}\tau<1$），也不满足判据。}

\item （推导）由泊松方程与玻尔兹曼因子 $n_e=n_0e^{e\phi/k_BT_e}$ 导出 $\lambda_D=\sqrt{\eps_0k_BT_e/(ne^2)}$。
\tkans{$\nabla^2\phi=-e(n_i-n_e)/\eps_0$；对 $|e\phi|\ll k_BT_e$ 展开 $n_e\approx n_0(1+e\phi/k_BT_e)$，代入得 $\nabla^2\phi=\phi/\lambda_D^2$，故 $\lambda_D=\sqrt{\eps_0k_BT_e/(n_0e^2)}$。工程式：$\lambda_D=7430\sqrt{T_e[\mathrm{eV}]/n[\mathrm{m^{-3}}]}\ \mathrm{m}$。}

\item （推导）证明德拜势在 $r=\lambda_D$ 处衰减到裸库仑势的 $1/e$，并求 $r=5\lambda_D$ 处剩余的百分比。
\tkans{$\phi\propto e^{-r/\lambda_D}/r$，故 $r=\lambda_D$ 时因子为 $e^{-1}=36.8\%$；$r=5\lambda_D$ 时为 $e^{-5}=0.67\%$，五个德拜长度之外屏蔽已经非常彻底，这就是「等离子体把电场关在自己家里」的定量含义。}

\item （估算）荧光灯正柱 $n=10^{16}\ \mathrm{m^{-3}}$、$kT_e=1\ \mathrm{eV}$：求 $\lambda_D$ 与 $N_D$，并判断它是否满足等离子体判据。
\tkans{$\lambda_D=7430\sqrt{1/10^{16}}=7.4\times10^{-5}\ \mathrm{m}$（74 $\mathrm{\mu m}$）；$N_D=\frac{4\pi}{3}\times10^{16}\times(7.43\times10^{-5})^3\approx1.7\times10^{4}\gg1$。管径量级为厘米，$\lambda_D\ll L$；电子温度 $1\ \mathrm{eV}$ 而气体接近室温，正是低温等离子体「电子热、离子冷」的典型参数。}

\item （估算）托卡马克边界（$n=10^{19}\ \mathrm{m^{-3}}$、$T_e=100\ \mathrm{eV}$）与日冕（$n=10^{15}\ \mathrm{m^{-3}}$、$T=100\ \mathrm{eV}$）的 $\lambda_D$ 之比是多少？
\tkans{$\lambda_D\propto\sqrt{T/n}$。边界：$\lambda_D=2.4\times10^{-5}\ \mathrm{m}$；日冕：$2.4\times10^{-3}\ \mathrm{m}$。比值 $\sqrt{10^{4}}=100$，即日冕的德拜长度比托卡马克边界大两个量级，密度低四个量级、温度相同。这解释了为什么日冕能容纳上百公里尺度的电流片。}
\end{enumerate}


\section{第 3 章：单粒子运动}
\tksrc{Chen 第 2 章、Fitzpatrick 第 2 章（平均化、绝热不变量）。}

\begin{enumerate}
\item （概念）$\vect{E}\times\vect{B}$ 漂移为什么与电荷符号和质量都无关？这会带来什么后果？
\tkans{运动方程 $m\diff\vect{v}/\diff t=q(\vect{E}+\vect{v}\times\vect{B})$ 的漂移解为 $\vect{v}_E=\vect{E}\times\vect{B}/B^2$，$q$ 与 $m$ 在求解过程中被消掉。后果：电子与离子\textbf{同向同速}漂移，不产生电荷分离，但整个等离子体一起平移，这是等离子体输运与「$E\times B$ 搅拌」的根源。}

\item （概念）磁镜为什么能反射粒子？损失锥的临界角由什么决定？
\tkans{磁矩 $\mu=mv_\perp^2/2B$ 是绝热不变量：$B$ 增大时 $v_\perp$ 增大，由能量守恒 $v_\parallel$ 减小，最终 $v_\parallel=0$ 而被反射。临界投掷角满足 $\sin^2\theta_c=B_{\min}/B_{\max}=1/R_m$（$R_m$ 为镜比），投掷角小于 $\theta_c$ 的粒子落入损失锥并逃逸。}

\item （推导）由 $m\diff\vect{v}/\diff t=q\vect{v}\times\vect{B}$（$\vect{B}=B\hat{z}$）导出 $\omega_c=|q|B/m$ 与 $\rho_L=v_\perp/\omega_c$，并说明回旋方向由什么决定。
\tkans{取 $\vect{v}=(v_\perp\cos\omega_ct,\,-v_\perp\sin\omega_ct,0)$ 代入即得圆周运动，角频率 $\omega_c=|q|B/m$，半径 $\rho_L=v_\perp/\omega_c=mv_\perp/|q|B$。旋转的\textbf{手性}（正电荷逆时针、负电荷顺时针，沿 $\vect{B}$ 方向看）由 $q$ 的符号决定，而频率与半径只依赖 $|q|/m$。}

\item （推导）托卡马克 $B=5\ \mathrm{T}$、氘离子 $T_i=10\ \mathrm{keV}$：求 $\omega_{ci}$、$v_i$ 与 $\rho_{Li}$。
\tkans{$\omega_{ci}=eB/m_i=1.602\times10^{-19}\times5/3.344\times10^{-27}=2.4\times10^{8}\ \mathrm{rad/s}$（$f_{ci}=38\ \mathrm{MHz}$）；$v_i=\sqrt{2k_BT_i/m_i}=9.8\times10^{5}\ \mathrm{m/s}$；$\rho_{Li}=v_i/\omega_{ci}=4.1\times10^{-3}\ \mathrm{m}$（4.1 mm）。离子回旋半径远小于装置半径（约 1--2 m），是磁约束可行的基本前提。}

\item （估算）地球磁层 $B=0.5\ \mathrm{G}$、$1\ \mathrm{keV}$ 质子：回旋半径是多少？为什么磁层能「困住」太阳风粒子？
\tkans{$\omega_{cp}=eB/m_p=1.602\times10^{-19}\times5\times10^{-5}/1.673\times10^{-27}=4.8\times10^{3}\ \mathrm{rad/s}$；$v=\sqrt{2k_BT/m_p}=4.4\times10^{5}\ \mathrm{m/s}$；$\rho_L=4.4\times10^{5}/4.8\times10^{3}=91\ \mathrm{m}$。它比磁层尺度（$10^{7}\ \mathrm{m}$）小五个量级，粒子只能沿磁力线螺旋前进，被磁场「串」在磁力线上，这就是辐射带与极光沉降的几何基础。}

\item （估算）$B=1\ \mathrm{T}$、$T_e=100\ \mathrm{eV}$ 的电子：回旋频率 $f_{ce}$ 与回旋半径 $\rho_{Le}$ 各是多少？与 $1\ \mathrm{cm}$ 探针尺度相比说明了什么？
\tkans{$\omega_{ce}=eB/m_e=1.76\times10^{11}\ \mathrm{rad/s}$，$f_{ce}=28\ \mathrm{GHz}$（毫米波段，可用于电子回旋加热与发射诊断）；$v_{te}=5.9\times10^{6}\ \mathrm{m/s}$，$\rho_{Le}=3.4\times10^{-5}\ \mathrm{m}$（34 $\mathrm{\mu m}$）。$\rho_{Le}\ll1\ \mathrm{cm}$：电子被强烈磁化，其横越磁力线的运动被冻结，只有沿磁力线才自由，这是磁化等离子体各向异性输运的起点。}
\end{enumerate}

\section{第 4 章：流体描述基础}
\tksrc{Chen 第 3 章、Fitzpatrick 第 4 章（Chapman--Enskog、Braginskii 输运）。}

\begin{enumerate}
\item （概念）连续性方程与动量方程各表达什么守恒？双流体模型还需要什么才能闭合？
\tkans{连续性方程 $\partial_tn+\nabla\cdot(n\vect{u})=0$ 是粒子数（质量）守恒；动量方程是动量守恒，含压强梯度、洛伦兹力 $\vect{j}\times\vect{B}$ 与碰撞摩擦。未知量多于方程，必须补\textbf{状态方程或能量方程}（等温 / 绝热假设）并加上麦克斯韦方程组，才能闭合。}

\item （概念）抗磁漂移的物理来源是什么？为什么它与电荷符号有关，而 $\vect{E}\times\vect{B}$ 漂移无关？
\tkans{压强梯度使回旋轨道两侧的受力不平衡，产生 $\vect{v}_D=-(\nabla p\times\vect{B})/(qnB^2)$：式中含 $q$，因此电子与离子\textbf{反向}漂移，形成抗磁电流（这正是等离子体电流与磁场的「自带」来源）。$\vect{E}\times\vect{B}$ 漂移不含 $q$，故与电荷符号无关。}

\item （推导）由流体动量方程导出抗磁漂移的标量形式 $v_{dia}=k_BT/(qBL_n)$（$L_n=|\nabla n/n|^{-1}$）。
\tkans{对稳态、无碰撞、各向同性压强，动量方程 $0=-\nabla p+qn(\vect{v}\times\vect{B})$ 给出 $\vect{v}_\perp=(\nabla p\times\vect{B})/(qnB^2)$。取 $\nabla p=k_BT\nabla n$、$|\nabla n|/n=1/L_n$，得 $v_{dia}=k_BT/(qBL_n)$；方向垂直于 $\nabla n$ 与 $\vect{B}$，电子与离子相反。}

\item （推导）$n=10^{19}\ \mathrm{m^{-3}}$、$T=100\ \mathrm{eV}$、$B=2\ \mathrm{T}$、$L_n=0.1\ \mathrm{m}$：求 $v_{dia}$，以及由此产生的抗磁电流密度 $j$。
\tkans{$v_{dia}=k_BT/(eBL_n)=1.602\times10^{-17}/(1.602\times10^{-19}\times2\times0.1)=500\ \mathrm{m/s}$。离子与电子速率相同、方向相反，故 $j=ne(v_i-v_e)=10^{19}\times1.602\times10^{-19}\times1000=1.6\times10^{3}\ \mathrm{A/m^2}$。1 cm 的密度梯度就能驱动千安每平方米量级的电流。}

\item （估算）氘等离子体 $T_e=10\ \mathrm{eV}$ 的离子声速 $c_s=\sqrt{k_BT_e/m_i}$ 是多少？与 $1\ \mathrm{eV}$ 的氩等离子体相比呢？
\tkans{氘（$m_i=3.34\times10^{-27}\ \mathrm{kg}$）：$c_s=\sqrt{T_e/m_i}=2.2\times10^{4}\ \mathrm{m/s}$；氩（$m_i=39.95\,m_p$）：$c_s=1.6\times10^{3}\ \mathrm{m/s}$。相差 14 倍，完全来自 $1/\sqrt{m_i}$，重离子等离子体的特征速度慢得多（鞘层、玻姆判据与预鞘加速都按这个速度标定）。}

\item （估算）若碰撞平均自由程 $\lambda_{mfp}=10^{-2}\ \mathrm{m}$、装置尺度 $L=0.5\ \mathrm{m}$，流体描述是否成立？
\tkans{$\lambda_{mfp}/L=0.02\ll1$：碰撞足够频繁，速度分布接近局域麦克斯韦分布，流体（局域热力学平衡）近似可用。反之若 $\lambda_{mfp}\gtrsim L$（如聚变芯部、空间等离子体），必须用动理学描述，这正是第 7 章与第 8 章的分工。}
\end{enumerate}


\section{第 5 章：磁流体力学}
\tksrc{Chen 第 3、6 章、Fitzpatrick 第 8 章（磁压、磁冻结、MHD 波），进阶参照 Freidberg \textit{Ideal MHD}。}

\begin{enumerate}
\item （概念）理想 MHD 的三条基本假设与两个「冻结」结论是什么？
\tkans{假设：准中性（$L\gg\lambda_D$）、单流体（离子与电子一起运动，$L\gg\rho_i$、$\omega\ll\omega_{ci}$）、理想欧姆定律 $\vect{E}+\vect{v}\times\vect{B}=0$（等价于 $R_m\gg1$）。结论：磁通量冻结（穿过任一流体回路的磁通守恒）、磁场与流体「焊」在一起；等价说法是磁力线像被流体拖着的橡皮筋。}

\item （概念）为什么说 $\vect{j}\times\vect{B}$ 可以写成「磁张力 + 磁压」？
\tkans{用矢量恒等式 $\vect{j}\times\vect{B}=(\vect{B}\cdot\nabla)\vect{B}/\mu_0-\nabla(B^2/2\mu_0)$。第一项沿磁力线方向，是\textbf{磁张力}（把弯曲的磁力线拉直）；第二项是各向同性的\textbf{磁压}（$B^2/2\mu_0$，把磁力线彼此推开）。于是 MHD 动量方程获得直观图像：磁力线是绷紧的弦，磁场是能推的气体。}

\item （推导）由感应方程 $\partial_t\vect{B}=\nabla\times(\vect{v}\times\vect{B})+(\eta/\mu_0)\nabla^2\vect{B}$ 导出磁扩散时间 $\tau_\eta=\mu_0L^2/\eta$ 与磁雷诺数 $R_m=\mu_0Lv/\eta$。
\tkans{忽略对流项，方程化为扩散方程，特征时间 $\tau_\eta=\mu_0L^2/\eta$；两项之比即 $R_m=\tau_\eta/\tau_{flow}=\mu_0Lv/\eta$。$R_m\gg1$ 时对流项主导，冻结成立；$R_m\ll1$ 时磁场自由扩散。}

\item （推导）$n=10^{20}\ \mathrm{m^{-3}}$、$B=5\ \mathrm{T}$ 的氘等离子体：求 Alfvén 速度，并用 $v_A/c$ 说明非相对论性。
\tkans{$v_A=B/\sqrt{\mu_0nm_i}=5/\sqrt{1.2566\times10^{-6}\times10^{20}\times3.344\times10^{-27}}=5/6.48\times10^{-7}=7.7\times10^{6}\ \mathrm{m/s}$；$v_A/c=2.6\%$。$v_A/c\ll1$ 只说明\textbf{非相对论 Alfv\'en 公式适用}（位移电流可以忽略），\textbf{不能据此断定磁冻结成立}：磁冻结要求理想欧姆定律成立，至少要 $R_m=\mu_0Lv/\eta\gg1$，还要检查霍尔项与电子压强项（正文第 5 章「MHD 适用条件：三道独立的门槛，别只看一道」）。本题未给温度、电阻率与流动尺度，无法判断（对聚变参数，$v_A$ 通常为 $10^{6}$--$10^{7}\ \mathrm{m/s}$）。}

\item （估算）$T_e=10\ \mathrm{eV}$、$\ln\Lambda=10$ 的氘等离子体：$\eta_\parallel\approx5.2\times10^{-5}\ln\Lambda\,T_e[\mathrm{eV}]^{-3/2}$，求 $\eta_\parallel$、$L=1\ \mathrm{m}$ 的 $\tau_\eta$，以及 $v=10^{4}\ \mathrm{m/s}$ 时的 $R_m$。磁冻结在这里成立吗？
\tkans{$\eta=5.2\times10^{-5}\times10/10^{1.5}=1.6\times10^{-5}\ \Omega\cdot\mathrm{m}$；$\tau_\eta=\mu_0L^2/\eta=1.2566\times10^{-6}/1.64\times10^{-5}=0.076\ \mathrm{s}$；$R_m=\mu_0Lv/\eta=1.2566\times10^{-6}\times1\times10^{4}/1.64\times10^{-5}=7.6\times10^{2}$。$R_m\gg1$ 成立，但只有两个多量级，若把 $L$ 缩到 1 cm 或把速度降到 $10^{3}\ \mathrm{m/s}$，$R_m$ 就跌到 10 以下，「理想」随即失效。结论：磁冻结是\textbf{需要检验}的近似，不是自动成立的定理。}

\item （估算）$\beta=2\mu_0nk_BT/B^2$：$n=10^{20}\ \mathrm{m^{-3}}$、$T=10\ \mathrm{keV}$、$B=5.3\ \mathrm{T}$ 时 $\beta$ 是多少？若要把 $\beta$ 压到 3\%（托卡马克运行区），需要多大磁场？
\tkans{$p=nk_BT=10^{20}\times10^{4}\times1.602\times10^{-19}=1.6\times10^{5}\ \mathrm{Pa}$；$\beta=2\times1.2566\times10^{-6}\times1.602\times10^{5}/5.3^2=1.4\%$。要 $\beta=3\%$：$B=\sqrt{2\mu_0p/0.03}=\sqrt{0.4026/0.03}=3.7\ \mathrm{T}$。可见同温度下提高磁场是压低 $\beta$（换取 MHD 稳定裕度）的直接手段。}
\end{enumerate}


\section{第 6 章：等离子体中的波}
\tksrc{Chen 第 4 章、Fitzpatrick 第 5、6、7 章（冷、非均匀、温等离子体），进阶参照 Stix \textit{Waves in Plasmas}。}

\begin{enumerate}
\item （概念）Bohm--Gross 色散 $\omega^2=\omega_{pe}^2+3k^2v_{te}^2$ 中两项各代表什么？$k\lambda_D\ll1$ 时回到什么？
\tkans{$\omega_{pe}^2$ 是电荷分离产生的静电恢复力（冷等离子体振荡）；$3k^2v_{te}^2$ 是电子热压强提供的恢复力（把 $3v_{te}^2$ 换成 $\gamma k_BT_e/m_e$ 即一维绝热声波）。$k\lambda_D\ll1$ 时热项可忽略，$\omega\to\omega_{pe}$：长波极限下电子集体振荡频率与 $k$ 无关。}

\item （概念）截止与共振有什么区别？为什么反射计能测量密度剖面？
\tkans{截止：$n$ 降到使 $k^2=0$（$\omega=\omega_{pe}$），波不能传播而被反射；共振：$k\to\infty$（如 $\omega=\omega_{ce}$ 附近的共振层），波被强烈吸收。反射计向等离子体发射扫频波，测到反射时的频率即给出截止层密度 $n_c(\omega)$；由不同频率的反射时间（或相位）反演出密度剖面。}

\item （推导）由 $\omega^2=\omega_{pe}^2+k^2c^2$ 导出截止密度 $n_c=\eps_0m_e\omega^2/e^2$，并求 $f=2.45\ \mathrm{GHz}$ 时的 $n_c$。
\tkans{令 $k=0$ 得 $\omega^2=\omega_{pe}^2=n_ce^2/(\eps_0m_e)$，故 $n_c=\eps_0m_e\omega^2/e^2$。$f=2.45\ \mathrm{GHz}$：$\omega=1.539\times10^{10}\ \mathrm{rad/s}$，$n_c=8.854\times10^{-12}\times9.109\times10^{-31}\times2.369\times10^{20}/2.566\times10^{-38}=7.4\times10^{16}\ \mathrm{m^{-3}}$。微波炉频率对应的正是这个量级的密度，等离子体工艺中「过密」与「欠密」的分界。}

\item （推导）离子声波 $\omega=kc_s$、$c_s=\sqrt{k_BT_e/m_i}$ 在什么条件下成立？求 $T_e=10\ \mathrm{eV}$ 时氘等离子体的 $c_s$。
\tkans{条件：$T_e\gg T_i$（离子被电子压强推着走）、$k\lambda_D\ll1$（准中性）。$c_s=\sqrt{T_e/m_i}=2.2\times10^{4}\ \mathrm{m/s}$，与第 4 章估算题一致；这个速度同时是玻姆速度与预鞘加速的标定速度。}

\item （估算）$n=10^{18}\ \mathrm{m^{-3}}$ 的等离子体频率 $f_{pe}$ 是多少？$2.45\ \mathrm{GHz}$ 的微波能穿透吗？
\tkans{$\omega_{pe}=\sqrt{ne^2/(\eps_0m_e)}=5.6\times10^{10}\ \mathrm{rad/s}$，$f_{pe}=9.0\ \mathrm{GHz}$。$f_{pe}>2.45\ \mathrm{GHz}$，即 $\omega<\omega_{pe}$，波是渐逝的，微波被反射（该等离子体对微波「过密」）。这正是「微波不能穿透高密度等离子体」的定量判据，也是 2.45 GHz 源难以维持高密度放电的原因。}

\item （估算）$n=10^{20}\ \mathrm{m^{-3}}$、$B=5\ \mathrm{T}$ 的氘等离子体中，Alfvén 波沿 $1\ \mathrm{m}$ 传播需要多久？
\tkans{$v_A=7.7\times10^{6}\ \mathrm{m/s}$（第 5 章已算），$t=L/v_A=1.3\times10^{-7}\ \mathrm{s}=0.13\ \mathrm{\mu s}$。微秒量级，比 MHD 不稳定性的增长时间快得多，说明理想 MHD 波可以在扰动增长期间多次往返，磁力线的「弹性」响应是瞬时的。}
\end{enumerate}

\section{第 7 章：动理学理论}
\tksrc{Chen 第 7 章、Fitzpatrick 第 3、7 章（Landau 阻尼、速度空间不稳定性），进阶参照 Krall \& Trivelpiece。}

\begin{enumerate}
\item （概念）弗拉索夫方程与玻尔兹曼方程差在哪里？为什么「无碰撞」也能有阻尼？
\tkans{弗拉索夫方程 $\partial_tf+\vect{v}\cdot\nabla f+(q/m)(\vect{E}+\vect{v}\times\vect{B})\cdot\nabla_vf=0$ 只保留自洽平均场，丢掉碰撞积分；玻尔兹曼方程多一项 $(\delta f/\delta t)_c$。Landau 阻尼不是碰撞阻尼，而是\textbf{波与共振粒子的可逆能量交换}：相速度附近的粒子被波加速或被减速，净效果是把相干波能转移为粒子的无规动能，因此熵增、波衰减，而无需任何碰撞。}

\item （概念）共振条件 $\omega-kv_\parallel=0$ 意味着什么？什么样的粒子主导 Landau 阻尼？
\tkans{平行速度等于波相速度 $v_\phi=\omega/k$ 的粒子与波「同相」共振，能持续交换能量。主导者是分布函数在 $v_\phi$ 处的斜率 $\partial f_0/\partial v$：斜率越平缓（或越靠尾部）贡献越小。$v_\phi\sim v_{te}$ 时阻尼最强；$v_\phi$ 远在分布尾部之外时阻尼指数小，这解释了为什么快波（$v_\phi\gg v_{te}$）几乎无阻尼。}

\item （推导）从弗拉索夫方程线性化（$f=f_0+\delta f$，$\delta f\propto e^{i(kx-\omega t)}$）导出 $\delta f=\dfrac{qE}{m}\dfrac{\partial f_0/\partial v}{i(kv-\omega)}$，并说明分母的极点如何处理。
\tkans{代入并保留一阶项得 $(-i\omega+ikv)\delta f+(qE/m)\partial_vf_0=0$，即上式。极点 $v=\omega/k$ 落在积分围道上，物理上对应共振粒子：绕开极点的方式（朗道围道，$\omega\to\omega+i0^+$）决定了积分取主值加半留数，其虚部给出阻尼/增长，极点的处理方式就是阻尼率的来源。}

\item （推导）写出麦克斯韦分布下的弱阻尼公式，并计算 $k\lambda_D=0.3$ 时的 $\gamma/\omega$。
\tkans{$\dfrac{\gamma}{\omega}=-\sqrt{\dfrac{\pi}{8}}\dfrac{1}{(k\lambda_D)^3}\exp\!\left[-\dfrac{1}{2k^2\lambda_D^2}-\dfrac{3}{2}\right]$。$k\lambda_D=0.3$：$(k\lambda_D)^{-3}=37.04$，指数 $\exp[-5.556-1.5]=e^{-7.056}=8.64\times10^{-4}$，故 $\gamma/\omega=-0.6266\times37.04\times8.64\times10^{-4}=-2.0\times10^{-2}$，即每周期衰减约 12\%。}

\item （估算）$k\lambda_D=0.25$ 与 $k\lambda_D=0.2$ 的 $\gamma/\omega$ 各是多少？为什么两者差近两个量级？
\tkans{$k\lambda_D=0.25$：$\gamma/\omega=-3.0\times10^{-3}$（每周期约 1.9\%）；$k\lambda_D=0.2$：$-6.5\times10^{-5}$（0.04\%）。差别几乎全部来自指数因子 $e^{-1/2k^2\lambda_D^2}$：从 $e^{-8}$ 降到 $e^{-12.5}$，倍率 $e^{-4.5}=0.011$。这就是「长波几乎无阻尼、$k\lambda_D\sim0.3$ 才明显」的定量含义，也是为什么必须把波长与 $\lambda_D$ 一起报告。}

\item （估算）双流不稳定性：$n=10^{18}\ \mathrm{m^{-3}}$、$\gamma_{\max}=\omega_{pe}/\sqrt{8}$，求 $\omega_{pe}$、$\gamma_{\max}$ 与特征增长时间。
\tkans{$\omega_{pe}=5.6\times10^{10}\ \mathrm{rad/s}$；$\gamma_{\max}=5.64\times10^{10}/2.828=2.0\times10^{10}\ \mathrm{s^{-1}}$；增长时间 $1/\gamma_{\max}=5.0\times10^{-11}\ \mathrm{s}=50\ \mathrm{ps}$。比任何碰撞或输运时间都快几个量级，这解释了为什么束-等离子体系统一旦满足条件就迅速进入强湍流（第 8 章「反常电阻」的动理学根源）。}
\end{enumerate}


\section{第 8 章：碰撞与输运}
\tksrc{Chen 第 5 章、Fitzpatrick 第 3、4 章（库仑对数、Rosenbluth 势），进阶参照 Helander \& Sigmar（新经典输运）。}

\begin{enumerate}
\item （概念）库仑对数 $\ln\Lambda$ 的物理意义是什么？为什么它几乎出现在所有碰撞率公式里？
\tkans{$\ln\Lambda=\ln(\lambda_D/b_{90})$，其中 $b_{90}$ 是使偏转角达 $90^\circ$ 的碰撞参数。库仑力长程，$90^\circ$ 散射的微分截面在小角发散，必须用\textbf{大量小角散射的累积}定义有效碰撞时间；积分给出对碰撞参数的对数权重 $\ln\Lambda$（典型 10--20，对 $n$、$T$ 只有弱依赖）。因此所有库仑碰撞率都带一个 $\ln\Lambda$ 因子，而它的不确定性（因子 2 以内）通常不影响结论。}

\item （概念）经典扩散与玻姆扩散的标度差异是什么？为什么实验上常观察到玻姆扩散？
\tkans{经典（碰撞）扩散 $D_\perp\propto n/T^{1/2}B^2$（对 $B$ 是 $1/B^2$），来自回旋运动被碰撞打断的随机行走；玻姆扩散 $D_B\propto k_BT_e/(eB)$（$1/B$），来自湍流涡的尺度 $\sim\rho_s$ 与寿命 $\sim1/\omega_*$。差异在 $B$ 的幂次：装置加大磁场时经典扩散迅速变小，而玻姆扩散只线性下降，因此在强磁场约束装置中湍流（玻姆类）输运常常主导。}

\item （推导）由 Spitzer 电阻率 $\eta_\parallel=5.2\times10^{-5}Z\ln\Lambda\,T_e[\mathrm{eV}]^{-3/2}\ \Omega\cdot\mathrm{m}$ 说明 $\eta\propto T^{-3/2}$，并计算温度从 $10\ \mathrm{eV}$ 升到 $1\ \mathrm{keV}$ 时 $\eta$ 下降多少倍。
\tkans{$\eta\propto T^{-3/2}$ 直接来自 $\eta\propto\nu_{ei}\propto nT^{-3/2}$（除以 $n$ 后仍留 $T^{-3/2}$）。倍率 $(1000/10)^{3/2}=100^{1.5}=1000$。温度升高 100 倍，电阻率降低 1000 倍，这是欧姆加热在高温下「自我熄火」的根源，也是为什么聚变装置必须靠中性束、射频等辅助加热。}

\item （推导）$n=10^{19}\ \mathrm{m^{-3}}$、$T_e=1\ \mathrm{keV}$、$\ln\Lambda=10$：求 $\eta_\parallel$ 与电子-离子碰撞时间 $\tau_{ei}$（$\nu_{ei}\approx2.9\times10^{-6}n\ln\Lambda\,T_e[\mathrm{eV}]^{-3/2}\ \mathrm{s^{-1}}$）。
\tkans{$\eta=5.2\times10^{-5}\times10/1000^{1.5}=1.6\times10^{-8}\ \Omega\cdot\mathrm{m}$；$\nu_{ei}=2.9\times10^{-6}\times10^{19}\times10/3.162\times10^{4}=9.2\times10^{9}\ \mathrm{s^{-1}}$，$\tau_{ei}=1/\nu_{ei}=1.1\times10^{-10}\ \mathrm{s}$。比较：$\omega_{pe}=5.6\times10^{9}\ \mathrm{rad/s}$，$\omega_{pe}\tau_{ei}=0.6$，量级为 1，说明该参数下集体效应与碰撞效应同等重要，「无碰撞」与「碰撞主导」都不是好的近似。}

\item （估算）上题电子的热速度 $v_{te}=1.9\times10^{7}\ \mathrm{m/s}$，平均自由程 $\lambda_{mfp}=v_{te}\tau_{ei}$ 是多少？与 $1\ \mathrm{m}$ 的装置尺度相比说明了什么？
\tkans{$\lambda_{mfp}=1.9\times10^{7}\times1.1\times10^{-10}=2.1\times10^{-3}\ \mathrm{m}$（2 mm），$\lambda_{mfp}/L=2\times10^{-3}\ll1$。虽然是「高温稀薄」等离子体，芯部的碰撞仍然足够频繁，局域输运系数（Spitzer、Braginskii）可用；但要注意这是\textbf{沿磁力线}的自由程，垂直方向的步长被回旋半径限制，两者相差 $\sim10^{4}$ 倍。}

\item （估算）中性粒子的电离/电荷交换截面取 $\sigma\approx10^{-19}\ \mathrm{m^2}$：边界 $n=10^{20}\ \mathrm{m^{-3}}$ 时穿透深度 $1/(n\sigma)$ 是多少？若边界密度升到 $10^{21}\ \mathrm{m^{-3}}$ 呢？
\tkans{$n=10^{20}$：$\lambda=1/(10^{20}\times10^{-19})=0.1\ \mathrm{m}$（10 cm）；$n=10^{21}$：$\lambda=1\ \mathrm{cm}$。穿透深度 $\propto1/n$，所以中性粒子（加料、再循环）只作用于边界薄层，芯部几乎完全电离，这也是偏滤器要把中性粒子「关」在靶板附近、以及中性束注入必须高能（才能边电离边穿入）的原因。}
\end{enumerate}


\section{第 9 章：等离子体不稳定性}
\tksrc{Chen 第 6、8 章、Fitzpatrick 第 7、9 章（撕裂模、重联），进阶参照 Freidberg 与 Wesson \textit{Tokamaks}。}

\begin{enumerate}
\item （概念）不稳定性的自由能来自哪里？「$\gamma>0$」的判据从哪里得到？
\tkans{自由能来源：密度/温度梯度（漂移波、交换模）、电流梯度（撕裂模、扭结模）、速度空间各向异性或束流（双流、损失锥不稳定性）、坏曲率与重力（交换模、瑞利--泰勒）。判据有两条路径：解色散关系看 $\omega$ 是否有正虚部（线性增长率），或用能量原理 $\delta W<0$（MHD 变分判据，不需要显式色散关系）。}

\item （概念）为什么「坏曲率」会驱动交换模？它与瑞利--泰勒不稳定性有什么类比？
\tkans{在曲率中心位于等离子体一侧（坏曲率）的区域，磁力线凹面朝向等离子体，磁压随半径下降得比等离子体压强快，等价于「重流体压在轻流体上」：把两根相邻磁通管交换位置会降低总能量，因此不稳定。类比关系是 $g_{eff}=c_s^2/R_c$ 扮演重力加速度，判据由 $\omega^2=-Akg_{eff}<0$ 给出，与流体瑞利--泰勒同形，差别只在「重力」来自磁力线弯曲而非真实引力。}

\item （推导）由 $\omega^2=-Akg$ 导出瑞利--泰勒增长率 $\gamma=\sqrt{Akg}$；再用 $g_{eff}=c_s^2/R_c$ 写出坏曲率交换模的增长率并说明稳定条件。
\tkans{$\omega^2<0$ 时 $\omega=\pm i\sqrt{Akg}$，故 $\gamma=\sqrt{Akg}$（$A$ 为 Atwood 数，$k$ 为波数）。交换模：$\gamma=\sqrt{kc_s^2/R_c}$。稳定条件：$k$ 与 $R_c$ 的组合使 $\gamma$ 被磁剪切或有限 $\beta$ 抑制，工程上表现为「坏曲率区足够小、磁剪切足够强」（把 $g_{eff}$ 变成沿磁力线的振荡平均，见磁井与平均最小 $B$）。}

\item （推导）漂移波色散 $\omega=\omega_*/(1+k_\perp^2\rho_s^2)$（$\omega_*=k_\perp\rho_sc_s/L_n$）：说明长波极限、增长率最大处的 $k_\perp$，以及湍流涡的尺度。
\tkans{$k_\perp\rho_s\ll1$ 时 $\omega\approx\omega_*$，频率正比于 $k_\perp$（沿抗磁方向传播）；$k_\perp\rho_s\gtrsim1$ 时离子有限回旋半径把涨落平均掉，$\omega$ 随 $k_\perp$ 下降，增长最强的尺度为 $k_\perp\rho_s\sim1$，即涡尺度 $\sim\rho_s$（离子声回旋半径，毫米量级）。这把「湍流输运」从定性图像落到一个可计算的尺度上。}

\item （估算）瑞利--泰勒：$A=1$、扰动波长 $6.3\ \mathrm{cm}$（$k=100\ \mathrm{m^{-1}}$）、$g=9.8\ \mathrm{m/s^2}$，求 $\gamma$ 与 e 折时间。
\tkans{$\gamma=\sqrt{Akg}=\sqrt{100\times9.8}=31\ \mathrm{s^{-1}}$，e 折时间 $1/\gamma=32\ \mathrm{ms}$。注意 $\gamma\propto\sqrt{k}$：波长越短的扰动长得越快（但受黏性/有限拉莫尔半径截断），这就是「界面扰动很快破碎成小尺度结构」的原因。}

\item （估算）托卡马克坏曲率区：$c_s=10^{5}\ \mathrm{m/s}$、$R_c=1\ \mathrm{m}$、$k=1\ \mathrm{m^{-1}}$，用 $\gamma=\sqrt{kc_s^2/R_c}$ 估算交换模增长率与增长时间。
\tkans{$g_{eff}=c_s^2/R_c=10^{10}\ \mathrm{m/s^2}$，$\gamma=\sqrt{kg_{eff}}=\sqrt{1.0\times10^{10}}=1.0\times10^{5}\ \mathrm{s^{-1}}$，增长时间 $10\ \mathrm{\mu s}$。比能量约束时间（$10^{-1}$--$10^{0}\ \mathrm{s}$）快四个量级：一旦坏曲率区与共振面重合，扰动会在远小于约束时间内长到非线性，所以托卡马克必须靠磁剪切与等离子体截面整形把「平均好曲率」做出来。}
\end{enumerate}

\section{第 10 章：约束与受控核聚变}
\tksrc{Chen 第 9、10 章、Fitzpatrick 第 8 章（太阳风、激波）；装置与工程量级参照 Goldston \& Rutherford 与 Wesson \textit{Tokamaks}。}

\begin{enumerate}
\item （概念）劳森判据的三个要素是什么？为什么写成「三乘积」而不是单一指标？
\tkans{密度 $n$、温度 $T$、能量约束时间 $\tau_E$。聚变功率密度 $\propto n^2\langle\sigma v\rangle$（两体过程），而单位体积热能 $\propto nT$、损失功率 $\propto nT/\tau_E$；两者之比要求 $nT\tau_E$ 足够大。三者的乘积形式说明\textbf{单靠提高密度或温度都不够}：还要把能量留住（$\tau_E$），这正是约束（磁场、惯性）而不是加热成为核心难点的原因。}

\item （概念）托卡马克与仿星器实现「旋转变换」的方式有何不同？各自的代价是什么？
\tkans{托卡马克靠\textbf{等离子体自身电流}产生极向场，从而让磁力线绕环螺旋（旋转变换）；仿星器完全用\textbf{外部三维线圈}的几何形状产生旋转变换。代价对比：托卡马克装置简单但需要驱动电流（射频/中性束），且电流剖面会驱动撕裂模、破裂等不稳定性；仿星器无净电流、无破裂风险，但线圈工程复杂、磁场误差敏感、输运优化困难。}

\item （推导）由零维功率平衡 $P_\alpha+P_{ext}=P_{loss}$ 与 $P_{fus}\propto n^2\langle\sigma v\rangle$ 说明点火条件如何化为 $nT\tau_E$ 的门槛，并给出 $Q$ 的定义。
\tkans{取 $P_\alpha=P_{fus}/5$（$\alpha$ 粒子带走 3.5 MeV/17.6 MeV），$P_{loss}=3nk_BT/\tau_E$。令 $P_{ext}=0$（点火）即得 $n^2\langle\sigma v\rangle E_\alpha/5=3nk_BT/\tau_E$，整理成 $nT\tau_E\ge\dfrac{15k_BT^2}{\langle\sigma v\rangle E_\alpha}$，右边只依赖 $T$，故门槛是三乘积。$Q=P_{fus}/P_{ext}$：$Q=1$ 得失相当，$Q=5$（ITER 目标）时 $\alpha$ 加热已占主导，$Q\to\infty$ 即点火。}

\item （推导）D--T 等离子体 $n_D=n_T=10^{20}\ \mathrm{m^{-3}}$、$\langle\sigma v\rangle=1.1\times10^{-22}\ \mathrm{m^3/s}$，每次反应释放 $17.6\ \mathrm{MeV}$：求聚变功率密度。
\tkans{$P_{fus}=n_Dn_T\langle\sigma v\rangle E_{fus}=10^{40}\times1.1\times10^{-22}\times(17.6\times10^{6}\times1.602\times10^{-19})=1.1\times10^{18}\times2.82\times10^{-12}=3.1\times10^{6}\ \mathrm{W/m^3}$，即约 $3\ \mathrm{MW/m^3}$。这是「堆芯量级」的功率密度，与裂变堆芯相当，这也是为什么聚变堆必须解决第一壁热流问题。}

\item （估算）接上题：$T=10\ \mathrm{keV}$ 时热能密度 $3nk_BT$ 是多少？$\alpha$ 加热功率密度 $P_\alpha$ 与点火门槛 $\tau_E$ 各是多少？
\tkans{$3nk_BT=3\times10^{20}\times10^{4}\times1.602\times10^{-19}=4.8\times10^{5}\ \mathrm{J/m^3}$；$P_\alpha=P_{fus}/5=6.2\times10^{5}\ \mathrm{W/m^3}$；自持要求 $\tau_E>3nk_BT/P_\alpha=0.77\ \mathrm{s}$。这是零维、忽略轫致辐射与杂质辐射的乐观估算；真实装置还要加上辐射损失与加热效率，门槛会更高（ITER 设计 $\tau_E\sim3\ \mathrm{s}$ 量级）。}

\item （估算）一个 $V=100\ \mathrm{m^3}$ 的 D--T 堆芯（参数同上）总聚变功率是多少？要达到 $1\ \mathrm{GW}$ 需要多大体积？
\tkans{$P=P_{fus}V=3.1\times10^{6}\times100=3.1\times10^{8}\ \mathrm{W}=310\ \mathrm{MW}$；要到 $1\ \mathrm{GW}$ 需 $V\approx320\ \mathrm{m^3}$（或把 $n$ 提高约 1.8 倍）。把「功率密度 × 体积」与工程尺度对上，是理解为什么聚变堆体积以千立方米计的最短路径。}
\end{enumerate}


\section{第 11 章：低温工业等离子体}
\tksrc{Chen 第 10 章、Fitzpatrick 第 4 章末（Langmuir 鞘层与探针），进阶参照 Lieberman \& Lichtenberg。}

\begin{enumerate}
\item （概念）低温等离子体为什么「电子热、离子冷」？这种非平衡状态带来了什么好处？
\tkans{电子质量小，在电场中被加速后与离子/中性粒子的弹性碰撞能量传递率正比于 $m_e/m_i$（$\sim10^{-4}$），因此电子从电场获得的能量难以传给重粒子：电子温度可达数个 eV，而气体温度接近室温。好处是\textbf{用低温衬底获得高能电子}，由它驱动电离、离解、激发与自由基生成，从而在材料表面实现刻蚀、沉积、改性（半导体工艺的基础）。}

\item （概念）鞘层为什么必然存在？浮动电位为什么是负的？
\tkans{等离子体与器壁接触时，电子热速度远大于离子玻姆速度（$v_{te}/v_B\sim\sqrt{m_i/m_e}\sim10^{2}$），初始电子通量远大于离子通量，器壁迅速充负电；负电位反过来排斥电子、吸引离子，直到两者通量相等。平衡电位即浮动电位，必然为负，量级为几个 $k_BT_e/e$。}

\item （推导）写出玻姆判据与 Child--Langmuir 定律，并给出由已知离子流密度求鞘层厚度的反解式。
\tkans{玻姆判据：离子进入鞘边时速度必须达到 $v_B=\sqrt{k_BT_e/m_i}$（否则鞘内离子密度下降快于电子，无法维持单调电位）。Child--Langmuir（无碰撞、空间电荷限制）$J_i=\dfrac{4\eps_0}{9}\sqrt{\dfrac{2e}{m_i}}\dfrac{|V_w|^{3/2}}{s^2}$，反解得 $s=\sqrt{\dfrac{4\eps_0}{9J_i}\sqrt{\dfrac{2e}{m_i}}|V_w|^{3/2}}$。鞘厚是 $\lambda_D$ 的若干倍，随 $|V_w|/T_e$ 的 $3/2$ 次方增长。}

\item （推导）氩等离子体 $T_e=3\ \mathrm{eV}$：求浮动电位 $V_f=\dfrac{k_BT_e}{2e}\ln\dfrac{m_i}{2\pi m_e}$ 与玻姆速度 $v_B$，并核对两者量级（$\sim5T_e/e$）。
\tkans{$\ln\dfrac{6.634\times10^{-26}}{2\pi\times9.109\times10^{-31}}=\ln(1.16\times10^{4})=9.36$，$V_f=1.5\times9.36=14.0\ \mathrm{V}$（$=4.7\,T_e/e$，符合 $\sim5T_e/e$ 的经验）。$v_B=\sqrt{3\times1.602\times10^{-19}/6.634\times10^{-26}}=2.7\times10^{3}\ \mathrm{m/s}$。}

\item （估算）$n_s=10^{16}\ \mathrm{m^{-3}}$、$T_e=3\ \mathrm{eV}$：求鞘边离子通量 $\Gamma=n_sv_B$ 与德拜长度，并说明为什么这个通量决定了刻蚀速率。
\tkans{$\Gamma=10^{16}\times2.69\times10^{3}=2.7\times10^{19}\ \mathrm{m^{-2}s^{-1}}$（每平方厘米每秒 $2.7\times10^{15}$ 个离子）；$\lambda_D=7430\sqrt{3/10^{16}}=1.3\times10^{-4}\ \mathrm{m}$。每个入射离子带来固定能量与反应概率，故刻蚀/沉积速率正比于 $\Gamma$；提高 $n$ 或 $T_e$ 都能提高通量，但 $T_e$ 同时决定离子能量（鞘电位降）。}

\item （估算）$n=10^{16}\ \mathrm{m^{-3}}$、$T_e=3\ \mathrm{eV}$：等离子体频率 $f_{pe}$ 与趋肤深度 $\delta=c/\omega_{pe}$ 各是多少？为什么 ICP 常用 $13.56\ \mathrm{MHz}$ 而 CCP 对频率不敏感？
\tkans{$\omega_{pe}=5.6\times10^{9}\ \mathrm{rad/s}$，$f_{pe}=9.0\times10^{8}\ \mathrm{Hz}=898\ \mathrm{MHz}$，$\delta=3\times10^{8}/5.6\times10^{9}=5.3\times10^{-2}\ \mathrm{m}$（5.3 cm）。$13.56\ \mathrm{MHz}\ll f_{pe}$，电磁波不能在等离子体中传播，只能以渐逝场（或碰撞/惯性趋肤层）耦合，ICP 靠线圈近场感应耦合，与放电室尺度、频率配合；CCP 靠鞘层电容耦合，离子能量主要由鞘电位降（而非频率）决定，因此对频率不敏感但加热效率受限。}
\end{enumerate}


\section{第 12 章：等离子体诊断}
\tksrc{Chen 第 10 章、Fitzpatrick 第 4 章；探针与波动诊断的严谨处理参照 Hutchinson \textit{Principles of Plasma Diagnostics}。}

\begin{enumerate}
\item （概念）单探针 $I$--$V$ 曲线的三段各对应什么物理过程？如何由曲线求 $T_e$ 与 $n_e$？
\tkans{离子饱和区（$V\ll V_p$）：几乎只收集离子，电流趋于 $I_{is}$；过渡区：电子被部分排斥，电流随 $V$ 指数上升；电子饱和区（$V\gg V_p$）：全部电子被收集，电流趋于 $I_{es}$。$T_e$ 由过渡区 $\ln I$ 对 $V$ 的斜率给出（斜率 $=e/k_BT_e$）；$n_e$ 由 $I_{es}=neA\bar{v}_e/4$（或由 $I_{is}$ 结合玻姆速度）反解。}

\item （概念）汤姆孙散射的 $\alpha$ 判据是什么？什么时候必须用集体散射理论？
\tkans{$\alpha=1/(k\lambda_D)$，$k$ 为散射波矢（$\approx$ 入射与散射波矢之差）。$\alpha\ll1$：散射发生在单个电子的尺度上，各电子独立散射，谱形直接反映电子速度分布（非集体/非相干散射）；$\alpha\gtrsim1$：散射体积内的德拜屏蔽与粒子关联不可忽略，谱中出现离子声波与等离子体波的特征线（集体散射），必须用带动态形状因子的理论解释。$\alpha\sim1$ 附近两种图像都不准，需完整理论。}

\item （推导）由双点法 $I(V_1)/I(V_2)=\exp[e(V_1-V_2)/k_BT_e]$ 给出 $T_e$ 的计算式，并说明取点原则。
\tkans{$T_e=\dfrac{e(V_1-V_2)}{k_B\ln[I(V_1)/I(V_2)]}$。取点须在\textbf{过渡区}且远离浮动电位与饱和区：太靠近 $V_f$ 时离子电流占比大，进入饱和后指数关系失效；同时两点电流比不宜过小（否则 $\ln$ 的相对误差被放大），也不宜跨越探针污染或二次电子发射明显的区域。}

\item （推导）某探针测得 $I(V)=50\ \mathrm{\mu A}$、$I(V+2\ \mathrm{V})=200\ \mathrm{\mu A}$：求 $T_e$；又知电子饱和电流 $I_{es}=1\ \mathrm{mA}$、收集面积 $A=1\ \mathrm{mm^2}$，估算 $n_e$。
\tkans{$T_e=2/\ln4=1.44\ \mathrm{eV}$。$v_{te}=\sqrt{2k_BT_e/m_e}=\sqrt{2\times1.443\times1.602\times10^{-19}/9.109\times10^{-31}}=7.1\times10^{5}\ \mathrm{m/s}$；$n_e=\dfrac{4I_{es}}{eA\bar{v}_e}$（以 $v_{te}$ 近似 $\bar v_e$）$=\dfrac{4\times10^{-3}}{1.602\times10^{-19}\times10^{-6}\times7.1\times10^{5}}=3.5\times10^{16}\ \mathrm{m^{-3}}$。注意用 $v_{te}$ 代替 $\bar v_e$ 会带来约 10\% 的系统偏差（严格用 $\bar v_e=\sqrt{8k_BT_e/\pi m_e}$）。}

\item （估算）微波干涉仪 $f=140\ \mathrm{GHz}$、$\bar n_e=10^{19}\ \mathrm{m^{-3}}$、路径长 $0.5\ \mathrm{m}$：求相移（弧度与条纹数）。
\tkans{截止密度 $n_c=\eps_0m_e\omega^2/e^2=2.43\times10^{20}\ \mathrm{m^{-3}}$；$\lambda=c/f=2.14\times10^{-3}\ \mathrm{m}$；相移 $\Delta\phi=\dfrac{\pi}{\lambda}\dfrac{\bar n_e}{n_c}L=\dfrac{\pi}{2.14\times10^{-3}}\times\dfrac{10^{19}}{2.43\times10^{20}}\times0.5=30.2\ \mathrm{rad}$，即 $30.2/2\pi=4.8$ 条条纹。条纹计数法能分辨约 0.1 条纹，对应的密度分辨约 $2\times10^{17}\ \mathrm{m^{-3}}$。}

\item （估算）磁探针：$N=10$ 匝、$A=10^{-4}\ \mathrm{m^2}$，磁场在 $100\ \mathrm{\mu s}$ 内从 0 升到 $1\ \mathrm{T}$：感应电压是多少？示波器（$\mathrm{mV}$ 灵敏度）能测吗？
\tkans{$\diff B/\diff t=10^{4}\ \mathrm{T/s}$；$V=NA\diff B/\diff t=10\times10^{-4}\times10^{4}=10\ \mathrm{V}$，远高于 $\mathrm{mV}$ 量级，甚至可能过载，实际探针常减小 $N$、$A$ 或加衰减。反向使用时，把 $V$ 数值积分即可还原 $B(t)$，这正是磁探针（及 Mirnov 线圈阵列）测磁涨落与 MHD 模的基本手段。}
\end{enumerate}


\section{来源、分级与版权说明}
\label{sec:app-l-source}

\begingroup\small
本附录的选题原则是\textbf{借鉴题型与分级，不复制题目}。下表逐章说明每题型的参照来源；表后是版权与独立性声明。

{\small
\begin{longtable}{@{}p{1.4cm}p{5.3cm}p{5.4cm}p{3.4cm}@{}}
\caption{分章题库的题型来源与难度参照（与附录 J 的逐章对照表 \ref{tab:app-j-map} 配套使用）}\label{tab:app-l-source}\\
\toprule
本书章 & 借鉴的教材与章节 & 借鉴的题型（不是题目） & 本题库构成 \\
\midrule
\endfirsthead
\toprule
本书章 & 借鉴的教材与章节 & 借鉴的题型（不是题目） & 本题库构成 \\
\midrule
\endhead
\bottomrule
\endlastfoot
1 & Chen 第 1 章及附录 A；Fitzpatrick 各章数学工具 & 分布函数积分、单位与量纲换算、屏蔽长度的量级题 & 概念 2 / 推导 2 / 估算 2 \\
2 & Chen 第 1 章；Fitzpatrick 第 1 章 & 判据的边界讨论、德拜球内粒子数估算 & 概念 2 / 推导 2 / 估算 2 \\
3 & Chen 第 2 章；Fitzpatrick 第 2 章 & 回旋频率与半径的数值题、损失锥几何 & 概念 2 / 推导 2 / 估算 2 \\
4 & Chen 第 3 章；Fitzpatrick 第 4 章 & 流体方程的守恒性讨论、抗磁漂移推导 & 概念 2 / 推导 2 / 估算 2 \\
5 & Chen 第 3、6 章；Fitzpatrick 第 8 章 & Alfvén 速度、磁扩散时间、$\beta$ 的标度题 & 概念 2 / 推导 2 / 估算 2 \\
6 & Chen 第 4 章；Fitzpatrick 第 5、6、7 章 & 截止密度、色散关系的极限讨论、波的传播时间 & 概念 2 / 推导 2 / 估算 2 \\
7 & Chen 第 7 章；Fitzpatrick 第 3、7 章 & 线性化推导、弱阻尼公式的数值代入与敏感度 & 概念 2 / 推导 2 / 估算 2 \\
8 & Chen 第 5 章；Fitzpatrick 第 3、4 章 & 库仑对数与碰撞时间、平均自由程与输运标度 & 概念 2 / 推导 2 / 估算 2 \\
9 & Chen 第 6、8 章；Fitzpatrick 第 7、9 章 & 增长率推导、瑞利--泰勒与交换模的类比题 & 概念 2 / 推导 2 / 估算 2 \\
10 & Chen 第 9、10 章；Goldston \& Rutherford；Wesson & 零维功率平衡、点火门槛、堆芯功率与体积估算 & 概念 2 / 推导 2 / 估算 2 \\
11 & Chen 第 10 章；Fitzpatrick 第 4 章末；Lieberman \& Lichtenberg & 鞘层与玻姆判据、浮动电位、通量与耦合方式 & 概念 2 / 推导 2 / 估算 2 \\
12 & Chen 第 10 章；Hutchinson & 探针特性曲线的读数、干涉仪相移、磁探针感应电压 & 概念 2 / 推导 2 / 估算 2 \\
\end{longtable}
}

\textbf{版权与独立性声明：}
\begin{itemize}
\item 上表所列「借鉴」仅指\textbf{题型分类与难度定位}（例如「求德拜长度的量级估算」「由玻姆判据讨论鞘层」这类公认的练习形式），属于学科共同的教学传统；
\item 本附录全部题目的\textbf{文字表述、参数选取与数值答案均由本书自拟}，并用仓库里的校核脚本逐条复算（两位有效数字）；
\item 本附录\textbf{不复制}任何教材的题目原文、解答、图表或数据表；引用教材仅给出书名与章节，便于读者回原书查找（完整书目见附录 E 与附录 J）；
\item 若读者发现某题与某教材的题目高度雷同，属独立命制的巧合，欢迎指出，我们会在下一版改写或补注出处。
\end{itemize}

\begin{tip}{做完 72 题之后}{tip-appl-next}
如果你把 12 章的 72 道题都做完了，可以回到附录 I 的第 5 节，把这 12 周的计划换成「每周一章：6 题自测 + 1 道附录 K 的范例精读」，题库与范例一一对应，正好覆盖全书主线。
\end{tip}
\endgroup


\end{document}
